% !TEX root = 第一章_数列极限(答案版).tex
\setcounter{chapter}{0}
\pagestyle{fancy}

\chapter{数列极限}


\section{预备知识}

作为培优教材,内容可能具有一定的综合性,先介绍可能会用到的几个公式:

\begin{proposition}[\marktwostar\ 对数不等式]
  当 $x>-1$ 时,
  \[
    \dfrac{x}{1+x}\leq \ln(1+x)\leq x.
  \]
  等号成立当且仅当 $x=0$ 时成立.
\end{proposition}

\begin{proposition}[\marktwostar]
  当 $0<x<\dfrac{\pi}{2}$ 时,
  \[
    \dfrac{2}{\pi}x\leq \sin x\leq x.
  \]
\end{proposition}

\begin{theorem}[\markstar\ Wallis 公式]
  \[
    \lim\limits_{n\to\infty}
    \dfrac{1}{2n+1}
    \left(\dfrac{(2n)!!}{(2n-1)!!}\right)^2
    =\dfrac{\pi}{2}.
  \]
\end{theorem}

\begin{theorem}[\markstar\ Stirling 公式]
  \[
    n!\sim \sqrt{2\pi n}\left(\dfrac{n}{e}\right)^n,
    \qquad n\to\infty.
  \]
  进一步,$n!$ 的估计还有中值形式:存在 $\theta_n\in(0,1)$,使得
  \[
    n!
    =\sqrt{2\pi n}\left(\dfrac{n}{e}\right)^n
    e^{\frac{\theta_n}{12n}}.
  \]
\end{theorem}

%例题1.1.1
\begin{example}
  证明命题 1.1.2.
\end{example}
\exampleblank

\begin{solution}
  命题 1.1.2 断言:当 $0<x<\dfrac{\pi}{2}$ 时,
  \[
    \dfrac{2}{\pi}x\leq \sin x\leq x.
  \]
  由于 $(\sin x)''=-\sin x<0$,故 $\sin x$ 在
  $\lbrack0,\dfrac{\pi}{2}\rbrack$ 上为凹函数.于是其图像位于连接
  $(0,0)$ 与 $(\dfrac{\pi}{2},1)$ 的弦的上方,从而
  \[
    \sin x\geq \dfrac{1-0}{\dfrac{\pi}{2}-0}x=\dfrac{2}{\pi}x.
  \]
  另一方面,由 Lagrange 中值定理,存在 $\xi\in(0,x)$,使得
  \[
    \sin x-\sin0=x\cos\xi\leq x.
  \]
  两式合并即得结论.
\end{solution}


\newpage

\section{收敛数列}


\subsection{证明极限存在}


\methodsection{(A) 结合拟合法,由定义证明}

首先介绍下拟合法,为了证明 $\displaystyle \lim\limits_{n\to\infty}x_n=A$,很多时候是转化成证明
$y_n=x_n-A$ 为无穷小(究其原因,主要是这样平移后一般更易放缩).故一般来说尽可能将
$x_n$ 的表达式化简.但有时 $x_n$ 虽然不能化简,$A$ 反而可以化简成类似 $x_n$ 的形式.

而对于 $\varepsilon-N$ 语言,我们一般有三种方法:

\begin{enumerate}
  \item (直接解不等式)对任意 $\varepsilon>0$,将 $|x_n-A|<\varepsilon$ 解出
        $n>N(\varepsilon)$,然后令 $N=N(\varepsilon)$ 即可;
  \item (放大法)有时 $|x_n-A|<\varepsilon$ 很难解出,则须将 $|x_n-A|$ 放缩化简成易于求解的
        $H(n)$,再解 $H(n)<\varepsilon$ 即可;
  \item (分段思想)有时 $|x_n-A|$ 特别复杂,我们须将其分为好几段(一般两到三段),再分别证每段为无穷小,而``段''的取法也是为了分段证无穷小的关键.
\end{enumerate}

%例题1.2.1
\begin{example}
  用 $\varepsilon-N$ 语言证明下列命题:
  \begin{enumerate}
    \item $\displaystyle \lim\limits_{n\to\infty}\sqrt[n]{n+1}=1$;
    \item 设 $\displaystyle \lim\limits_{n\to\infty}x_n=a$,则
          $\displaystyle \lim\limits_{n\to\infty}\sqrt[3]{x_n}=\sqrt[3]{a}$;
    \item 设 $\displaystyle \lim\limits_{n\to\infty}x_n=A$,则
          \[
            \lim\limits_{n\to\infty}\dfrac{x_1+x_2+\cdots+x_n}{n}=A.
          \]
  \end{enumerate}
\end{example}
\exampleblank

\begin{solution}
  依次证明三问.

  (1) 显然 $\sqrt[n]{n+1}>1$.任给 $\varepsilon>0$,当 $n\geq2$ 时,由二项式展开有
  \[
    (1+\varepsilon)^n
    \geq C_n^2\varepsilon^2
    =\dfrac{n(n-1)}{2}\varepsilon^2.
  \]
  取 $N$ 充分大,使 $n>N$ 时
  $\dfrac{n(n-1)}{2}\varepsilon^2>n+1$,则
  \[
    1<\sqrt[n]{n+1}<1+\varepsilon.
  \]
  故 $\displaystyle \lim\limits_{n\to\infty}\sqrt[n]{n+1}=1$.

  (2) 若 $a\neq0$,则 $x_n$ 最终与 $a$ 同号,因而
  \[
    \left|\sqrt[3]{x_n}-\sqrt[3]{a}\right|
    =\dfrac{|x_n-a|}
    {x_n^{\frac{2}{3}}+\sqrt[3]{x_n}\sqrt[3]{a}+a^{\frac{2}{3}}}
    \leq\dfrac{|x_n-a|}{|a|^{\frac{2}{3}}}.
  \]
  给定 $\varepsilon>0$,取 $n$ 充分大使右端小于 $\varepsilon$,即得结论.
  若 $a=0$,则对充分大的 $n$ 有 $|x_n|<\varepsilon^3$,故
  $|\sqrt[3]{x_n}|<\varepsilon$.所以结论对一切 $a\in\mathbb{R}$ 均成立.

  (3) 任给 $\varepsilon>0$,取 $N_0$,使得 $k>N_0$ 时
  $|x_k-A|<\dfrac{\varepsilon}{2}$.令
  \[
    M=\sum\limits_{k=1}^{N_0}|x_k-A|.
  \]
  当 $n>N_0$ 时,
  \[
    \left|\dfrac{1}{n}\sum\limits_{k=1}^n x_k-A\right|
    \leq \dfrac{M}{n}+\dfrac{1}{n}\sum\limits_{k=N_0+1}^n|x_k-A|
    <\dfrac{M}{n}+\dfrac{\varepsilon}{2}.
  \]
  再取 $n>\dfrac{2M}{\varepsilon}$,便得所求差小于 $\varepsilon$,故算术平均数收敛于 $A$.
\end{solution}

%例题1.2.2
\begin{example}
  设 $\displaystyle \lim\limits_{n\to\infty}a_n=a$,
  \[
    x_n=\dfrac{a_1+2a_2+\cdots+na_n}{1+2+\cdots+n},
  \]
  则 $\displaystyle \lim\limits_{n\to\infty}x_n=a$.
\end{example}
\exampleblank

\begin{solution}
  记 $P_n=1+2+\cdots+n=\dfrac{n(n+1)}{2}$.任给 $\varepsilon>0$,取
  $N_0$,使 $k>N_0$ 时 $|a_k-a|<\dfrac{\varepsilon}{2}$.再令
  \[
    C=\sum\limits_{k=1}^{N_0}k|a_k-a|.
  \]
  于是当 $n>N_0$ 时,
  \[
    |x_n-a|
    \leq \dfrac{C}{P_n}
    +\dfrac{1}{P_n}\sum\limits_{k=N_0+1}^n k|a_k-a|
    <\dfrac{C}{P_n}+\dfrac{\varepsilon}{2}.
  \]
  由于 $P_n\to+\infty$,上式第一项最终小于 $\dfrac{\varepsilon}{2}$,故
  $|x_n-a|<\varepsilon$.因此 $x_n\to a$.
\end{solution}

%例题1.2.3
\begin{example}
  证明:若 $p_k>0\ (k=1,2,\ldots)$,且
  \[
    \lim\limits_{n\to\infty}\dfrac{p_n}{p_1+p_2+\cdots+p_n}=0,
    \qquad
    \lim\limits_{n\to\infty}a_n=a,
  \]
  则
  \[
    \lim\limits_{n\to\infty}
    \dfrac{p_1a_n+p_2a_{n-1}+\cdots+p_na_1}{p_1+p_2+\cdots+p_n}=a.
  \]
\end{example}
\exampleblank

\begin{solution}
  记 $P_n=\displaystyle \sum\limits_{k=1}^n p_k$,并把所求表达式改写为
  \[
    T_n=\dfrac{1}{P_n}\sum\limits_{j=1}^n p_{n+1-j}a_j.
  \]
  先注意到,对每个固定的 $j$,当 $n\to\infty$ 时 $n+1-j\to\infty$,从而
  \[
    0\leq\dfrac{p_{n+1-j}}{P_n}
    \leq\dfrac{p_{n+1-j}}{P_{n+1-j}}\longrightarrow0.
  \]
  任给 $\varepsilon>0$,取 $N_0$,使 $j>N_0$ 时 $|a_j-a|<\varepsilon$.
  于是
  \[
    |T_n-a|
    \leq \sum\limits_{j=1}^{N_0}
    \dfrac{p_{n+1-j}}{P_n}|a_j-a|
    +\varepsilon\sum\limits_{j=N_0+1}^{n}\dfrac{p_{n+1-j}}{P_n}.
  \]
  右端有限项之和趋于 $0$,后一项不超过 $\varepsilon$.因此
  $\varlimsup\limits_{n\to\infty}|T_n-a|\leq\varepsilon$.由 $\varepsilon$ 的任意性知 $T_n\to a$.
\end{solution}

%例题1.2.4
\begin{example}
  设 $x\to0$ 时 $f(x)\sim x$,
  \[
    x_n=\sum\limits_{i=1}^n f\left(\dfrac{2i-1}{n^2}a\right),
  \]
  试证:
  \[
    \lim\limits_{n\to\infty}x_n=a\qquad(a>0).
  \]
\end{example}
\exampleblank

\begin{solution}
  任给 $\varepsilon\in(0,1)$.由 $\dfrac{f(x)}{x}\to1$,存在 $\delta>0$,使得
  $0<x<\delta$ 时
  \[
    (1-\varepsilon)x<f(x)<(1+\varepsilon)x.
  \]
  因为
  \[
    0<\dfrac{(2i-1)a}{n^2}
    \leq\dfrac{(2n-1)a}{n^2}\longrightarrow0
  \]
  对 $1\leq i\leq n$ 一致成立,所以当 $n$ 充分大时
  \[
    (1-\varepsilon)\dfrac{(2i-1)a}{n^2}
    <f\left(\dfrac{(2i-1)a}{n^2}\right)
    <(1+\varepsilon)\dfrac{(2i-1)a}{n^2}.
  \]
  对 $i=1,\ldots,n$ 求和,并利用
  $\displaystyle\sum\limits_{i=1}^n(2i-1)=n^2$,得到
  \[
    (1-\varepsilon)a<x_n<(1+\varepsilon)a.
  \]
  由 $\varepsilon$ 的任意性和夹逼定理可知 $x_n\to a$.
\end{solution}

\begin{remark}
  当 $x\to0$ 时,$\sin x,\tan x,\arcsin x,\arctan x,e^x-1,\ln(1+x)$ 都与 $x$ 等价,因此 $f$ 取这些函数时结论都成立,如
  \[
    \lim\limits_{n\to\infty}\sum\limits_{i=1}^n
    \sin\dfrac{2i-1}{n^2}a=a.
  \]
\end{remark}

%例题1.2.5
\begin{example}
  设 $\displaystyle \lim\limits_{n\to\infty}a_n=a$,试证:
  \[
    \lim\limits_{n\to\infty}\dfrac{1}{2^n}
    \left(
    C_{n}^{0}a_0+C_{n}^{1}a_1+\cdots+C_{n}^{k}a_k+\cdots+C_{n}^{n}a_n
    \right)=a.
  \]
\end{example}
\exampleblank

\begin{solution}
  记
  \[
    T_n=\dfrac{1}{2^n}\sum\limits_{k=0}^n C_n^k a_k.
  \]
  任给 $\varepsilon>0$,取 $N_0$,使 $k>N_0$ 时 $|a_k-a|<\varepsilon$.
  则
  \[
    |T_n-a|
    \leq\dfrac{1}{2^n}\sum\limits_{k=0}^{N_0}C_n^k|a_k-a|
    +\varepsilon\dfrac{1}{2^n}\sum\limits_{k=N_0+1}^n C_n^k.
  \]
  第二项不超过 $\varepsilon$.对每个固定的 $k$,
  $C_n^k$ 是关于 $n$ 的 $k$ 次多项式,故 $\dfrac{C_n^k}{2^n}\to0$;因此第一项也趋于 $0$.
  于是 $\varlimsup\limits_{n\to\infty}|T_n-a|\leq\varepsilon$,从而 $T_n\to a$.
\end{solution}

%例题1.2.6
\begin{example}
  设函数 $f(x),g(x)$ 在 $0$ 的某个邻域里有定义,$g(x)>0$,
  \[
    \lim\limits_{x\to0}\dfrac{f(x)}{g(x)}=1,
  \]
  且 $n\to\infty$ 时,$\alpha_{mn}\rightrightarrows0\ (m=1,2,\ldots,n)$,亦即对任意
  $\varepsilon>0$,存在 $N(\varepsilon)>0$,当 $n>N(\varepsilon)$ 时,对一切
  $m=1,2,\ldots,n$,有 $|\alpha_{mn}|<\varepsilon$.另设 $\alpha_{mn}\neq0$,试证:
  \[
    \lim\limits_{n\to\infty}\sum\limits_{m=1}^n f(\alpha_{mn})
    =
    \lim\limits_{n\to\infty}\sum\limits_{m=1}^n g(\alpha_{mn}),
  \]
  当右端极限存在时成立.
\end{example}
\exampleblank

\begin{solution}
  由 $\dfrac{f(x)}{g(x)}\to1$,任给 $\varepsilon\in(0,1)$,存在 $\delta>0$,使得
  $0<|x|<\delta$ 时
  \[
    (1-\varepsilon)g(x)<f(x)<(1+\varepsilon)g(x).
  \]
  由 $\alpha_{mn}\rightrightarrows0$,当 $n$ 充分大时,对所有
  $1\leq m\leq n$ 均有 $0<|\alpha_{mn}|<\delta$,又 $g(\alpha_{mn})>0$.因此
  \[
    (1-\varepsilon)\sum\limits_{m=1}^n g(\alpha_{mn})
    <\sum\limits_{m=1}^n f(\alpha_{mn})
    <(1+\varepsilon)\sum\limits_{m=1}^n g(\alpha_{mn}).
  \]
  若右端原极限为有限值 $L$,则 $L\geq0$.令 $n\to\infty$,取下,上极限得
  \[
    (1-\varepsilon)L
    \leq\varliminf\limits_{n\to\infty}\sum\limits_{m=1}^n f(\alpha_{mn})
    \leq\varlimsup\limits_{n\to\infty}\sum\limits_{m=1}^n f(\alpha_{mn})
    \leq(1+\varepsilon)L.
  \]
  再令 $\varepsilon\to0^+$,即知左端和式的极限存在且等于 $L$.
\end{solution}

\begin{remark}
\[
\dfrac{\sqrt[3]{1+x}-1}{\dfrac{x}{3}}         
=\dfrac{3}{(1+x)^{\frac{2}{3}}+(1+x)^{\frac{1}{3}}+1}
\longrightarrow 1 \qquad (x\to 0).
\]

  在本题中取
  \[
    f(x)=\sqrt[3]{1+x}-1,
    \qquad
    g(x)=\dfrac{x}{3},
    \qquad
    \alpha_{i,n}=\dfrac{i}{n^2}.
  \]
  因 $\displaystyle\max_{1\leq i\leq n}\alpha_{i,n}=\dfrac{1}{n}\to0$,由上面的结论,
  \[
    \lim\limits_{n\to\infty}\sum\limits_{i=1}^n
    \left(\sqrt[3]{1+\dfrac{i}{n^2}}-1\right)
    =
    \lim\limits_{n\to\infty}\sum\limits_{i=1}^n\dfrac{i}{3n^2}
    =\dfrac{1}{6}.
  \]
\end{remark}



\methodsection{(B) Cauchy 收敛准则}

再回顾一下 Cauchy 收敛准则:
\[
  \{x_n\}\text{ 收敛}
  \iff
  \forall\varepsilon>0,\ \exists N>0,
  \quad m,n>N\Longrightarrow |x_n-x_m|<\varepsilon,
\]
等价地,
\[
  \{x_n\}\text{ 收敛}
  \iff
  \forall\varepsilon>0,\ \exists N>0,
  \quad n>N\Longrightarrow |x_{n+p}-x_n|<\varepsilon\quad(\forall p>0).
\]

\begin{remark}
  值得注意的是,后者的 $p$ 应是关于 $p\in\mathbb{N}$ 一致成立,即
  \[
    \{x_n\}\text{ 收敛}
    \iff
    |x_{n+p}-x_n|\rightrightarrows0\qquad(n\to\infty).
  \]
  关于 $p$ 一致成立,此时的 $N$ 与 $p$ 无关.
\end{remark}

由于 Cauchy 收敛准则的形式,在判别和式和不易求得极限值的数列敛散性时尤为有效.

%例题1.2.7
\begin{example}
  判断如下命题真伪:
  \[
    \text{数列 }\{a_n\}\text{ 存在极限 }\lim\limits_{n\to\infty}a_n=a
    \iff
    \forall p\in\mathbb{N},\ \lim\limits_{n\to\infty}|a_{n+p}-a_n|=0.
  \]
\end{example}
\exampleblank

\begin{solution}
  命题的充分性错误,必要性正确.

  若 $a_n\to a$,则对任意固定的 $p\in\mathbb{N}$,
  \[
    |a_{n+p}-a_n|
    \leq|a_{n+p}-a|+|a_n-a|\longrightarrow0.
  \]
  反之不成立.取 $a_n=\ln n$,则对任意固定的 $p$,
  \[
    |a_{n+p}-a_n|=\ln\left(1+\dfrac{p}{n}\right)\longrightarrow0,
  \]
  但 $a_n\to+\infty$,故 $\{a_n\}$ 不收敛.问题在于 Cauchy 准则要求对所有
  $m,n$ 同时一致控制,而题设只逐个固定了间隔 $p$.
\end{solution}

%例题1.2.8
\begin{example}
  设
  \[
    x_n=\dfrac{\sin1}{2}+\dfrac{\sin2}{2^2}+\cdots+\dfrac{\sin n}{2^n},
  \]
  试证 $\{x_n\}$ 收敛.
\end{example}
\exampleblank

\begin{solution}
  当 $m>n$ 时,
  \[
    |x_m-x_n|
    \leq\sum\limits_{k=n+1}^m\dfrac{|\sin k|}{2^k}
    \leq\sum\limits_{k=n+1}^{\infty}\dfrac{1}{2^k}
    =\dfrac{1}{2^n}.
  \]
  右端与 $m$ 无关且趋于 $0$,故 $\{x_n\}$ 满足 Cauchy 收敛准则,从而收敛.
\end{solution}

%例题1.2.9
\begin{example}
  设
  \[
    x_n=\sum\limits_{k=2}^n\dfrac{\cos k}{k(k-1)},
  \]
  试证 $\{x_n\}$ 收敛.
\end{example}
\exampleblank

\begin{solution}
  当 $m>n\geq2$ 时,
  \[
    |x_m-x_n|
    \leq\sum\limits_{k=n+1}^m\dfrac{1}{k(k-1)}
    =\sum\limits_{k=n+1}^m\left(\dfrac{1}{k-1}-\dfrac{1}{k}\right)
    =\dfrac{1}{n}-\dfrac{1}{m}<\dfrac{1}{n}.
  \]
  故 $\{x_n\}$ 为 Cauchy 数列,所以收敛.
\end{solution}



\methodsection{(C) Toeplitz 定理}

\begin{theorem}[\markstar\ Toeplitz 定理]
  设对 $n,k\in\mathbb{Z}^+$,有 $t_{nk}\geq0$,又
  \[
    \sum\limits_{k=1}^n t_{nk}=1,
    \qquad
    \lim\limits_{n\to\infty}t_{nk}=0.
  \]
  若已知 $\displaystyle \lim\limits_{n\to\infty}a_n=a$,定义
  \[
    x_n=\sum\limits_{k=1}^n t_{nk}a_k,\qquad n\in\mathbb{Z}^+,
  \]
  则
  \[
    \lim\limits_{n\to\infty}x_n=a.
  \]
\end{theorem}

\begin{remark}
  该定理还有几种类型:
  \begin{enumerate}
    \item 将 $\displaystyle \sum\limits_{k=1}^n t_{nk}=1$ 改为
          $\displaystyle \lim\limits_{n\to\infty}\sum\limits_{k=1}^n t_{nk}=1$;
    \item 不要求 $t_{nk}$ 非负,将 (1) 中的条件改为:存在 $M>0$,使得对任意 $n$,有
          \[
            |t_{n1}|+|t_{n2}|+\cdots+|t_{nn}|\leq M.
          \]
  \end{enumerate}
\end{remark}

%例题1.2.10
\begin{example}
  由 Toeplitz 定理导出 Stolz 定理.
\end{example}
\exampleblank

\begin{solution}
  设 $\{y_n\}$ 严格递增且 $y_n\to+\infty$,并设
  \[
    \dfrac{x_n-x_{n-1}}{y_n-y_{n-1}}\longrightarrow A.
  \]
  令
  \[
    d_k=y_k-y_{k-1}>0,
    \qquad
    c_k=\dfrac{x_k-x_{k-1}}{d_k}.
  \]
  则
  \[
    \dfrac{x_n-x_0}{y_n-y_0}
    =\dfrac{\sum\limits_{k=1}^n d_kc_k}{\sum\limits_{k=1}^n d_k}.
  \]
  右端是 $c_1,\ldots,c_n$ 的正权平均,且权和趋于无穷.由 Toeplitz 定理,该式趋于 $A$.
  又因 $y_n\to+\infty$,
  \[
    \dfrac{x_n}{y_n}-\dfrac{x_n-x_0}{y_n-y_0}
    =\dfrac{x_0y_n-x_ny_0}{y_n(y_n-y_0)}\longrightarrow0,
  \]
  其中最后一步可由前式有界推出 $x_n=O(y_n)$.故
  $\dfrac{x_n}{y_n}\to A$,这就是 Stolz 定理的无穷型.
\end{solution}

\begin{remark}
  因此许多 Stolz 定理能求的问题,Toeplitz 也能,甚至还能解决前面定义里分段思想解决的问题.
\end{remark}

%例题1.2.11
\begin{example}
  设
  \[
    A_n=\sum\limits_{k=1}^n a_k,
  \]
  $\{A_n\}$ 收敛,又正数列 $\{p_n\}$ 单调递增,且为无穷大量,证明:
  \[
    \lim\limits_{n\to\infty}
    \dfrac{p_1a_1+p_2a_2+\cdots+p_na_n}{p_n}=0.
  \]
\end{example}
\exampleblank

\begin{solution}
  由分部求和公式,
  \[
    \sum\limits_{k=1}^n p_ka_k
    =p_nA_n-\sum\limits_{k=1}^{n-1}(p_{k+1}-p_k)A_k.
  \]
  将上式改写为
  \[
    \dfrac{1}{p_n}\sum\limits_{k=1}^n p_ka_k
    =A_n-\sum\limits_{k=1}^{n-1}
    \dfrac{p_{k+1}-p_k}{p_n}A_k.
  \]
  各权非负,其和为 $\dfrac{p_n-p_1}{p_n}\to1$,且对固定 $k$ 权趋于 $0$.
  由 Toeplitz 定理,右端第二项与 $A_n$ 有相同极限,故两者之差趋于 $0$.
\end{solution}

%例题1.2.12
\begin{example}
  设 $0<\lambda<1$,$\{a_n\}$ 收敛于 $a$,证明:
  \[
    \lim\limits_{n\to\infty}
    \left(a_n+\lambda a_{n-1}+\lambda^2a_{n-2}+\cdots+\lambda^n a_0\right)
    =\dfrac{a}{1-\lambda}.
  \]
\end{example}
\exampleblank

\begin{solution}
  将原式乘以趋于 $1-\lambda$ 的因子
  $\dfrac{1-\lambda}{1-\lambda^{n+1}}$,得到
  \[
    \dfrac{1-\lambda}{1-\lambda^{n+1}}
    \sum\limits_{j=0}^n\lambda^j a_{n-j}
    =\sum\limits_{k=0}^n
    \dfrac{(1-\lambda)\lambda^{n-k}}{1-\lambda^{n+1}}a_k.
  \]
  右端各权非负,权和为 $1$;对每个固定的 $k$,
  \[
    \dfrac{(1-\lambda)\lambda^{n-k}}{1-\lambda^{n+1}}
    \longrightarrow0.
  \]
  由 Toeplitz 定理,上式右端趋于 $a$.又
  $\dfrac{1-\lambda}{1-\lambda^{n+1}}\to1-\lambda$,故题中和式趋于
  $\dfrac{a}{1-\lambda}$.
\end{solution}

%例题1.2.13
\begin{example}
  设 $\displaystyle \lim\limits_{n\to\infty}x_n=0$,存在 $k$,使得
  \[
    |y_1|+|y_2|+\cdots+|y_n|\leq k.
  \]
  对任意 $n$ 成立,令
  \[
    z_n=x_1y_n+x_2y_{n-1}+\cdots+x_ny_1.
  \]
  证明:$\displaystyle \lim\limits_{n\to\infty}z_n=0$.
\end{example}
\exampleblank

\begin{solution}
  将 $z_n$ 写成
  \[
    z_n=\sum\limits_{j=1}^n t_{nj}x_j,
    \qquad t_{nj}=y_{n+1-j}.
  \]
  对每个固定的 $j$,由
  $\displaystyle\sum\limits_{m=1}^{\infty}|y_m|<+\infty$ 可知
  $t_{nj}=y_{n+1-j}\to0$;并且
  \[
    \sum\limits_{j=1}^n|t_{nj}|
    =\sum\limits_{j=1}^n|y_j|\leq k.
  \]
  Toeplitz 定理的零极限型只需固定列趋于 $0$ 且各行绝对值之和一致有界.
  因 $x_j\to0$,立即得到 $z_n\to0$.
\end{solution}



\methodsection{(D) 单调有界原理}

%例题1.2.14
\begin{example}
  证明:数列
  \[
    x_n=1+\dfrac{1}{2}+\dfrac{1}{3}+\cdots+\dfrac{1}{n}-\ln n
    \qquad(n=1,2,\ldots).
  \]
  单调递减有界,从而有极限.
\end{example}
\exampleblank

\begin{solution}
  由 $\ln(1+t)>\dfrac{t}{1+t}$ 对 $t>0$ 成立,
  \[
    x_{n+1}-x_n
    =\dfrac{1}{n+1}-\ln\left(1+\dfrac{1}{n}\right)<0,
  \]
  故 $\{x_n\}$ 单调递减.另一方面,由于 $\dfrac{1}{x}$ 在 $(0,+\infty)$ 上递减,
  \[
    \ln n=\int_1^n\dfrac{1}{x}\,\mathrm{d}x
    \leq\sum\limits_{k=1}^{n-1}\dfrac{1}{k},
  \]
  从而 $x_n\geq\dfrac{1}{n}>0$.所以 $\{x_n\}$ 单调递减且有下界,必有有限极限.
\end{solution}

%例题1.2.15
\begin{example}
  设
  \[
    x_n=\sum\limits_{k=1}^n\dfrac{1}{\sqrt{k}}-2\sqrt{n},
  \]
  证明:$\displaystyle \lim\limits_{n\to\infty}x_n$ 存在.
\end{example}
\exampleblank

\begin{solution}
  先证单调性.由
  \[
    2(\sqrt{n+1}-\sqrt n)
    =\dfrac{2}{\sqrt{n+1}+\sqrt n}
    >\dfrac{1}{\sqrt{n+1}},
  \]
  有
  \[
    x_{n+1}-x_n
    =\dfrac{1}{\sqrt{n+1}}-2(\sqrt{n+1}-\sqrt n)<0.
  \]
  再由 $x^{-\frac{1}{2}}$ 递减,
  \[
    \sum\limits_{k=1}^n\dfrac{1}{\sqrt k}
    \geq\int_1^{n+1}\dfrac{1}{\sqrt x}\,\mathrm{d}x
    =2\sqrt{n+1}-2.
  \]
  故
  \[
    x_n\geq2\sqrt{n+1}-2-2\sqrt n>-2.
  \]
  于是 $\{x_n\}$ 单调递减且有下界,从而极限存在.
\end{solution}



\methodsection{(E) 利用上,下极限 *}

这部分方法由于预备知识理论性较强,故放在本章最后一节作为选学内容.



\subsection{证明极限不存在}

证明极限不存在,最本源的当然是由定义判断,但事实上,由过程繁琐等原因,反而用的较少.



\methodsection{(A) 无界数列发散}

%例题1.2.16
\begin{example}
  证明:
  \[
    S_n=1+\dfrac{1}{2}+\dfrac{1}{3}+\cdots+\dfrac{1}{n}.
  \]
  无界,进而 $\{S_n\}$ 发散.
\end{example}
\exampleblank

\begin{solution}
  取 $n=2^m$.把调和和按二进区间分组,得
  \[
    S_{2^m}
    =1+\sum\limits_{j=1}^m
    \sum\limits_{k=2^{j-1}+1}^{2^j}\dfrac{1}{k}
    \geq1+\sum\limits_{j=1}^m2^{j-1}\dfrac{1}{2^j}
    =1+\dfrac{m}{2}.
  \]
  故子列 $\{S_{2^m}\}$ 无界,从而 $\{S_n\}$ 无界.收敛数列必有界,所以
  $\{S_n\}$ 发散.
\end{solution}



\methodsection{(B) 数列与子列的关系}

(1) 存在发散子列;
(2) 存在不收敛于同一极限的两个子列.

%例题1.2.17
\begin{example}
  若 $\{a_n\}$ 中三个子列 $\{a_{2n}\},\{a_{2n-1}\},\{a_{3n}\}$ 皆收敛,则 $\{a_n\}$ 收敛.
\end{example}
\exampleblank

\begin{solution}
  设
  \[
    a_{2n}\to\alpha,
    \qquad a_{2n-1}\to\beta,
    \qquad a_{3n}\to\gamma.
  \]
  子列 $\{a_{6n}\}$ 既是 $\{a_{2n}\}$ 的子列,又是 $\{a_{3n}\}$ 的子列,故
  $\alpha=\gamma$.子列 $\{a_{6n-3}\}$ 既是 $\{a_{2n-1}\}$ 的子列,又是
  $\{a_{3n}\}$ 的子列,故 $\beta=\gamma$.于是 $\alpha=\beta$.
  原数列的偶数项与奇数项均趋于同一极限,所以 $a_n\to\gamma$.
\end{solution}

%例题1.2.18
\begin{example}
  设
  \[
    a_n=\dfrac{1}{n}-\dfrac{2}{n}+\dfrac{3}{n}+\cdots+(-1)^{n-1}\dfrac{n}{n}
    \qquad(n=1,2,\ldots),
  \]
  则 $\{a_n\}$ 发散.
\end{example}
\exampleblank

\begin{solution}
  当 $n=2m$ 时,
  \[
    a_{2m}=\dfrac{(1-2)+(3-4)+\cdots+((2m-1)-2m)}{2m}
    =-\dfrac{1}{2}.
  \]
  当 $n=2m-1$ 时,
  \[
    a_{2m-1}
    =\dfrac{(1-2)+\cdots+((2m-3)-(2m-2))+(2m-1)}{2m-1}
    =\dfrac{m}{2m-1}\longrightarrow\dfrac{1}{2}.
  \]
  偶子列与奇子列极限不同,故 $\{a_n\}$ 发散.
\end{solution}

%例题1.2.19
\begin{example}
  证明:$\{\sin n\}$ 发散.
\end{example}
\exampleblank

\begin{solution}

  \begin{theorem}[Dirichlet 逼近定理]
    设 $\alpha$ 为无理数,则存在无穷多个互不相同的有理数 $\dfrac{p}{q}$
    ($p\in\mathbb{Z}$,$q\in\mathbb{Z}^+$),使得
    \[
      \left|\alpha-\dfrac{p}{q}\right|<\dfrac{1}{q^2}.
    \]
  \end{theorem}
  

  取 $\alpha=\dfrac{1}{2\pi}$,由上述定理可选出 $q$ 严格递增的一列
  $\dfrac{p_j}{q_j}$,则 $\dfrac{1}{q_j}\to0$,并且
  \[
    \left|\dfrac{q_j}{2\pi}-p_j\right|
    =q_j\left|\dfrac{1}{2\pi}-\dfrac{p_j}{q_j}\right|
    <\dfrac{1}{q_j}\longrightarrow0.
  \]
  即 $q_j-2\pi p_j\to0$.由 $\sin$ 以 $2\pi$ 为周期,
  \[
    \sin q_j=\sin(q_j-2\pi p_j)\longrightarrow0.
  \]
  另一方面,
  \[
    \sin(q_j+1)\longrightarrow\sin1\neq0.
  \]
  因此 $\{\sin n\}$ 有两个极限不同的子列,所以它不收敛.
\end{solution}



\methodsection{(C) 利用恒等式反证得结论}

该方法多用于证关于三角函数列的发散问题.

%例题1.2.20
\begin{example}
  证明:$\{\sin n\}$ 发散.
\end{example}
\exampleblank

\begin{solution}
  反设 $\sin n\to A$.由恒等式
  \[
    \sin(n+2)-2\cos1\sin(n+1)+\sin n=0
  \]
  令 $n\to\infty$,得到 $2A(1-\cos1)=0$,故 $A=0$.
  再由
  \[
    \sin(n+1)-\cos1\sin n=\sin1\cos n
  \]
  知 $\cos n\to0$.这与恒等式 $\sin^2n+\cos^2n=1$ 矛盾.
  故 $\{\sin n\}$ 发散.
\end{solution}

%例题1.2.21
\begin{example}
  证明:$\{\tan n\}$ 发散.
\end{example}
\exampleblank

\begin{solution}
  反设 $\tan n\to A$.记 $t=\tan1\neq0$.由正切加法公式,
  \[
    \tan(n+1)=\dfrac{\tan n+t}{1-t\tan n}.
  \]
  若 $1-tA=0$,则上式分母趋于 $0$,而分子不可能同时趋于 $0$,这与
  $\tan(n+1)$ 有有限极限矛盾.故 $1-tA\neq0$,令 $n\to\infty$ 得
  \[
    A=\dfrac{A+t}{1-tA}.
  \]
  整理得 $t(A^2+1)=0$,不可能成立.所以 $\{\tan n\}$ 发散.
\end{solution}

%例题1.2.22
\begin{example}
  若存在极限
  \[
    \lim\limits_{n\to\infty}(a\sin n+b\cos n)=A,
  \]
  则 $a=b=0$.
\end{example}
\exampleblank

\begin{solution}
  记
  \[
    u_n=a\sin n+b\cos n.
  \]
  若 $u_n\to A$,由
  \[
    u_{n+2}-2\cos1\,u_{n+1}+u_n=0
  \]
  可得 $A=0$.又令
  \[
    v_n=a\cos n-b\sin n,
  \]
  则
  \[
    u_{n+1}-\cos1\,u_n=\sin1\,v_n.
  \]
  所以 $v_n\to0$.然而
  \[
    u_n^2+v_n^2=a^2+b^2,
  \]
  令 $n\to\infty$ 得 $a^2+b^2=0$.因此 $a=b=0$.
\end{solution}



\methodsection{(D) Cauchy 收敛准则}

%例题1.2.23
\begin{example}
  证明:数列
  \[
    S_n=\sum\limits_{k=1}^n\dfrac{1}{k}\qquad(n=1,2,\ldots).
  \]
  发散.
\end{example}
\exampleblank

\begin{solution}
  对任意 $n$,
  \[
    S_{2n}-S_n=\sum\limits_{k=n+1}^{2n}\dfrac{1}{k}
    \geq n\dfrac{1}{2n}=\dfrac{1}{2}.
  \]
  因此对 $\varepsilon=\dfrac{1}{4}$,无论 $N$ 多大,取 $n>N$,$m=2n$,总有
  $|S_m-S_n|\geq\dfrac{1}{2}>\varepsilon$.故 $\{S_n\}$ 不满足 Cauchy 准则,从而发散.
\end{solution}

%例题1.2.24
\begin{example}
  设
  \[
    x_n=1+\sqrt{\dfrac{1}{2}}+\sqrt{\dfrac{1}{3}}+\cdots+\sqrt{\dfrac{1}{n}},
  \]
  求证:$\{x_n\}$ 发散.
\end{example}
\exampleblank

\begin{solution}
  当 $m=2n$ 时,
  \[
    x_{2n}-x_n
    =\sum\limits_{k=n+1}^{2n}\dfrac{1}{\sqrt k}
    \geq\dfrac{n}{\sqrt{2n}}=\sqrt{\dfrac{n}{2}}.
  \]
  该差不但不趋于 $0$,反而趋于 $+\infty$,故 $\{x_n\}$ 不满足 Cauchy 准则,
  所以发散.
\end{solution}



\subsection{求极限值}


\methodsection{(A) 迫敛性}

%例题1.2.25
\begin{example}
  求下列数列 $\{a_n\}$ 的极限 $\displaystyle \lim\limits_{n\to\infty}a_n$:
  \begin{enumerate}
    \item $\displaystyle a_n=\sqrt[n]{2\sin^2n+\cos^2n}$;
    \item $\displaystyle a_n=(n+1+n\cos n)^{\frac{1}{2n+n\sin n}}$.
  \end{enumerate}
\end{example}
\exampleblank

\begin{solution}
  (1) 因为 $1\leq2\sin^2n+\cos^2n\leq2$,所以
  \[
    1\leq\sqrt[n]{2\sin^2n+\cos^2n}\leq\sqrt[n]{2}\longrightarrow1.
  \]
  由夹逼定理,所求极限为 $1$.

  (2) 由于
  \[
    1\leq n+1+n\cos n\leq2n+1,
    \qquad
    n\leq2n+n\sin n\leq3n,
  \]
  故
  \[
    0\leq
    \ln\left[(n+1+n\cos n)^{\frac{1}{2n+n\sin n}}\right]
    =\dfrac{\ln(n+1+n\cos n)}{2n+n\sin n}
    \leq\dfrac{\ln(2n+1)}{n}\longrightarrow0.
  \]
  从而原式趋于 $e^0=1$.
\end{solution}

%例题1.2.26
\begin{example}
  求下列数列 $\{a_n\}$ 的极限 $\displaystyle \lim\limits_{n\to\infty}a_n$:
  \begin{enumerate}
    \item $\displaystyle a_n=\sqrt[n]{1+b^n}\quad(b>0)$;
    \item $\displaystyle a_n=\sqrt[n]{1+b^n+\left(\dfrac{b^2}{2}\right)^n}\quad(b>0)$.
  \end{enumerate}
\end{example}
\exampleblank

\begin{solution}
  (1) 若 $0<b\leq1$,则
  \[
    1\leq\sqrt[n]{1+b^n}\leq\sqrt[n]{2}\longrightarrow1.
  \]
  若 $b>1$,则
  \[
    b\leq\sqrt[n]{1+b^n}\leq b\sqrt[n]{2}\longrightarrow b.
  \]
  故极限为 $\max\{1,b\}$.

  (2) 令 $M=\max\left\{1,b,\dfrac{b^2}{2}\right\}$.和式中的三项均不超过 $M^n$,且至少一项等于
  $M^n$,所以
  \[
    M\leq\sqrt[n]{1+b^n+\left(\dfrac{b^2}{2}\right)^n}
    \leq\sqrt[n]{3}\,M.
  \]
  故极限为 $M$.
\end{solution}

%例题1.2.27
\begin{example}[\markstar]
  设
  \[
    A=\max\{a_1,a_2,\ldots,a_m\},\qquad a_k>0,
  \]
  求
  \[
    \lim\limits_{n\to\infty}\sqrt[n]{a_1^n+a_2^n+\cdots+a_m^n}.
  \]
\end{example}
\exampleblank

\begin{solution}
  由 $A=\max\{a_1,a_2,\ldots,a_m\}$ 及各 $a_j>0$,有
  \[
    A^n\leq a_1^n+a_2^n+\cdots+a_m^n\leq mA^n.
  \]
  开 $n$ 次方得
  \[
    A\leq\sqrt[n]{a_1^n+a_2^n+\cdots+a_m^n}\leq\sqrt[n]{m}\,A.
  \]
  由 $\sqrt[n]{m}\to1$ 及夹逼定理,所求极限为 $A$.
\end{solution}

%例题1.2.28
\begin{example}[\markstar]
  设 $f(x)>0$ 在 $\lbrack 0,1\rbrack$ 上连续,则
  \[
    \lim\limits_{n\to\infty}
    \sqrt[n]{\dfrac{1}{n}\sum\limits_{i=1}^n\left[f\left(\dfrac{i}{n}\right)\right]^n}
    =
    \max_{0\leq x\leq1}f(x).
  \]
\end{example}
\exampleblank

\begin{solution}

  由连续性,$f$ 在 $\lbrack0,1\rbrack$ 上取得最大值.记

  \[
    M=\max\limits_{0\leq x\leq1}f(x)>0.
  \]

  显然

  \[
    0<\left[\dfrac{1}{n}\sum\limits_{k=1}^n f^n\left(\dfrac{k}{n}\right)\right]^{\frac{1}{n}}
    \leq M.
  \]

  任给 $\varepsilon\in(0,M)$,在某个长度为 $\delta>0$ 的区间上有
  $f(x)>M-\varepsilon$.当 $n$ 充分大时,该区间内至少有一个分点
  $\dfrac{k_n}{n}$,故

  \[
    \left[\dfrac{1}{n}\sum\limits_{k=1}^n f^n\left(\dfrac{k}{n}\right)\right]^{\frac{1}{n}}
    \geq\dfrac{1}{n^{\frac{1}{n}}}f\left(\dfrac{k_n}{n}\right)
    >\dfrac{M-\varepsilon}{n^{\frac{1}{n}}}.
  \]

  由 $n^{\frac{1}{n}}\to1$,取下极限得其不小于 $M-\varepsilon$.再令
  $\varepsilon\to0^+$,即得极限为 $M$.

\end{solution}

%例题1.2.29
\begin{example}[\markstar]
  设 $\{a_n\}$ 是正实数列,
  \[
    a=\sup\{a_1,a_2,\ldots\},
  \]
  求证:
  \[
    \lim\limits_{n\to\infty}\left(\sum\limits_{k=1}^n a_k^n\right)^{\frac{1}{n}}=a.
  \]
\end{example}
\exampleblank

\begin{solution}
  先设 $a<+\infty$.因为 $0<a_k\leq a$,有
  \[
    \sqrt[n]{a_1^n+\cdots+a_n^n}\leq\sqrt[n]{na^n}=\sqrt[n]{n}\,a\longrightarrow a.
  \]
  另一方面,对任意固定的 $k$,当 $n\geq k$ 时
  \[
    \sqrt[n]{a_1^n+\cdots+a_n^n}\geq a_k.
  \]
  故其下极限不小于 $\sup_k a_k=a$,从而极限为 $a$.

  若 $a=+\infty$,则对任意 $M>0$ 可取固定的 $k$ 使 $a_k>M$;上式说明所求数列的下极限不小于
  $M$.因 $M$ 任意,此时极限为 $+\infty$.
\end{solution}

%例题1.2.30
\begin{example}
  求下列数列 $\{a_n\}$ 的极限 $\displaystyle \lim\limits_{n\to\infty}a_n$:
  \begin{enumerate}
    \item $\displaystyle a_n=\dfrac{1}{n^n}\sum\limits_{k=1}^n k^k$;
    \item $\displaystyle a_n=\dfrac{\sqrt[n]{\sum\limits_{k=1}^n k^n}}{n}$;
    \item $\displaystyle a_n=\sum\limits_{k=1}^n\dfrac{k}{n^2+n+k}$.
  \end{enumerate}
\end{example}
\exampleblank

\begin{solution}
  (1) 下界由和式最后一项给出:
  \[
    \dfrac{1}{n^n}\sum\limits_{k=1}^n k^k\geq1.
  \]
  又因 $k^k\leq n^k$,
  \[
    \dfrac{1}{n^n}\sum\limits_{k=1}^n k^k
    \leq\sum\limits_{k=1}^n n^{k-n}
    =\sum\limits_{j=0}^{n-1}\dfrac{1}{n^j}
    \leq\dfrac{1}{1-\dfrac{1}{n}}\longrightarrow1.
  \]
  故极限为 $1$.

  (2) 有
  \[
    n^n\leq\sum\limits_{k=1}^n k^n\leq n\cdot n^n.
  \]
  因此
  \[
    1\leq\dfrac{1}{n}\sqrt[n]{\sum\limits_{k=1}^n k^n}\leq\sqrt[n]{n}\longrightarrow1.
  \]

  (3) 对 $1\leq k\leq n$,
  \[
    n^2+n+1\leq n^2+n+k\leq n^2+2n.
  \]
  所以
  \[
    \dfrac{n(n+1)}{2(n^2+2n)}
    \leq\sum\limits_{k=1}^n\dfrac{k}{n^2+n+k}
    \leq\dfrac{n(n+1)}{2(n^2+n+1)}.
  \]
  两端均趋于 $\dfrac{1}{2}$,故所求极限为 $\dfrac{1}{2}$.
\end{solution}



\methodsection{(B) 利用积分定义}

%例题1.2.31
\begin{example}
  求下列极限:
  \begin{enumerate}
    \item $\displaystyle \lim\limits_{n\to\infty}\sum\limits_{k=1}^n\dfrac{1}{n+k}$;
    \item $\displaystyle \lim\limits_{n\to\infty}\sum\limits_{k=1}^n
            \dfrac{\sin\dfrac{k\pi}{n}}{n+\dfrac{k}{n}}$.
  \end{enumerate}
\end{example}
\exampleblank

\begin{solution}
  (1) 写成 Riemann 和:
  \[
    \sum\limits_{k=1}^n\dfrac{1}{n+k}
    =\dfrac{1}{n}\sum\limits_{k=1}^n\dfrac{1}{1+\dfrac{k}{n}}
    \longrightarrow\int_0^1\dfrac{1}{1+x}\,\mathrm{d}x=\ln2.
  \]

  (2) 同理,
  \[
    \sum\limits_{k=1}^n\dfrac{\sin\dfrac{k\pi}{n}}{n+\dfrac{k}{n}}
    =\dfrac{1}{n}\sum\limits_{k=1}^n
    \dfrac{\sin\dfrac{k\pi}{n}}{1+\dfrac{k}{n^2}}.
  \]
  因 $\dfrac{1}{1+\dfrac{k}{n^2}}\to1$ 对 $1\leq k\leq n$ 一致成立,故上式与
  $\displaystyle \dfrac{1}{n}\sum\limits_{k=1}^n\sin\dfrac{k\pi}{n}$ 之差趋于 $0$.因此极限为
  \[
    \int_0^1\sin(\pi x)\,\mathrm{d}x=\dfrac{2}{\pi}.
  \]
\end{solution}

%例题1.2.32
\begin{example}
  求下列极限:
  \begin{enumerate}
    \item $\displaystyle
            \lim\limits_{n\to\infty}
            \dfrac{\sqrt[n]{(n+1)(n+2)\cdots(n+n)}}{n}$;
    \item $\displaystyle
            \lim\limits_{n\to\infty}
            \dfrac{\sqrt[n]{n(n+1)\cdots(2n-1)}}{n}$;
    \item $\displaystyle
            \lim\limits_{n\to\infty}\dfrac{\sqrt[n]{n!}}{n}$.
  \end{enumerate}
\end{example}
\exampleblank

\begin{solution}
  (1) 取对数后得到
  \[
    \dfrac{1}{n}\sum\limits_{k=1}^n\ln\left(1+\dfrac{k}{n}\right)
    \longrightarrow\int_0^1\ln(1+x)\,\mathrm{d}x=2\ln2-1.
  \]
  故原极限为 $e^{2\ln2-1}=\dfrac{4}{e}$.

  (2) 取对数后为
  \[
    \dfrac{1}{n}\sum\limits_{k=0}^{n-1}\ln\left(1+\dfrac{k}{n}\right),
  \]
  它有相同的积分极限 $2\ln2-1$,故原极限仍为 $\dfrac{4}{e}$.

  (3) 取对数后为
  \[
    \dfrac{1}{n}\ln\dfrac{n!}{n^n}
    =\dfrac{1}{n}\sum\limits_{k=1}^n\ln\dfrac{k}{n}.
  \]
  这是广义积分 $\displaystyle \int_0^1\ln x\,\mathrm{d}x=-1$ 的 Riemann 和;也可由 Stirling 公式直接得到其极限为 $-1$.
  故原极限为 $e^{-1}$.
\end{solution}



\methodsection{(C) 利用级数性质}

主要有两个运用: (1) 收敛级数的通项趋于零;(2) 
$\sum\limits_{n=1}^{\infty}(x_{n+1}-x_{n})$
绝对收敛$\Rightarrow $
$\displaystyle \lim\limits_{n\to\infty}x_n$
存在.

%例题1.2.33
\begin{example}
  求下列极限 $\displaystyle \lim\limits_{n\to\infty}x_n$:
  \begin{enumerate}
    \item $\displaystyle x_n=\dfrac{5^n\cdot n!}{(2n)^n}$;
    \item $\displaystyle x_n=\dfrac{11\cdot12\cdot13\cdots(n+10)}{2\cdot5\cdot8\cdots(3n-1)}$.
  \end{enumerate}
\end{example}
\exampleblank

\begin{solution}
  由 Stirling 公式 $n!\sim\sqrt{2\pi n}\left(\dfrac{n}{e}\right)^n$,有
  \[
    x_n\sim\sqrt{2\pi n}\left(\dfrac{5}{2e}\right)^n\longrightarrow0,
  \]
  因为 $\dfrac{5}{2e}<1$.

  (2) 题中乘积的相邻两项之比满足
  \[
    \dfrac{x_{n+1}}{x_n}=\dfrac{n+11}{3n+2}\longrightarrow\dfrac{1}{3}<1.
  \]
  取 $q\in\left(\dfrac{1}{3},1\right)$.当 $n$ 充分大时 $x_{n+1}\leq qx_n$,故
  $x_n\leq Cq^n\to0$.
\end{solution}

%例题1.2.34
\begin{example}
  设
  \[
    x_n=\dfrac{1}{2\ln2}+\cdots+\dfrac{1}{n\ln n}-\ln\ln n,
  \]
  求证:$\{x_n\}$ 收敛.
\end{example}
\exampleblank

\begin{solution}
  令 $f(x)=\dfrac{1}{x\ln x}$.当 $x>e$ 时 $f$ 单调递减.
  当 $n$ 充分大时,
  \[
    x_{n+1}-x_n
    =f(n+1)-\int_n^{n+1}f(x)\,\mathrm{d}x\leq0,
  \]
  故 $\{x_n\}$ 最终单调递减.另一方面,
  \[
    \sum\limits_{k=2}^n f(k)
    \geq\int_2^{n+1}f(x)\,\mathrm{d}x
    =\ln\ln(n+1)-\ln\ln2,
  \]
  所以
  \[
    x_n\geq\ln\dfrac{\ln(n+1)}{\ln n}-\ln\ln2,
  \]
  右端有统一下界.因此 $\{x_n\}$ 收敛,所给极限存在.
\end{solution}

%例题1.2.35
\begin{example}
  设
  \[
    0\leq x_{n+1}\leq x_n+\dfrac{1}{n^2},
  \]
  则 $\{x_n\}$ 收敛.
\end{example}
\exampleblank

\begin{solution}
  令
  \[
    y_n=x_n+\sum\limits_{k=n}^{\infty}\dfrac{1}{k^2}.
  \]
  由题设,
  \[
    y_{n+1}
    =x_{n+1}+\sum\limits_{k=n+1}^{\infty}\dfrac{1}{k^2}
    \leq x_n+\dfrac{1}{n^2}+\sum\limits_{k=n+1}^{\infty}\dfrac{1}{k^2}
    =y_n.
  \]
  又 $y_n\geq0$,故 $\{y_n\}$ 单调递减且有下界,从而存在极限 $L\geq0$.
  由于 $\displaystyle \sum\limits_{k=n}^{\infty}\dfrac{1}{k^2}\to0$,有
  \[
    x_n=y_n-\sum\limits_{k=n}^{\infty}\dfrac{1}{k^2}\longrightarrow L.
  \]
  故 $\{x_n\}$ 收敛.
\end{solution}



\methodsection{(D) Stolz 定理}

Stolz 可以视为数列版本的洛必达法则,在解决和式及 $\dfrac{0}{0}$,$\dfrac{*}{\infty}$ 未定式极限问题十分有效.

\begin{theorem}[\markstar\ Stolz 定理]
  \leavevmode
  \begin{enumerate}
    \item ($\dfrac{*}{\infty}$ 型)设 $\{y_n\}$ 严格单调递增,且
          $\displaystyle \lim\limits_{n\to\infty}y_n=+\infty$.若
          \[
            \lim\limits_{n\to\infty}
            \dfrac{x_n-x_{n-1}}{y_n-y_{n-1}}=A
            \qquad(A\text{ 为有限数或 }\pm\infty),
          \]
          则 $\displaystyle \lim\limits_{n\to\infty}\dfrac{x_n}{y_n}=A$.
    \item ($\dfrac{0}{0}$ 型)设 $\{y_n\}$ 严格单调递减,且
          $\displaystyle \lim\limits_{n\to\infty}x_n
          =\lim\limits_{n\to\infty}y_n=0$.若
          \[
            \lim\limits_{n\to\infty}
            \dfrac{x_n-x_{n-1}}{y_n-y_{n-1}}=A
            \qquad(A\text{ 为有限数}),
          \]
          则 $\displaystyle \lim\limits_{n\to\infty}\dfrac{x_n}{y_n}=A$.
  \end{enumerate}
\end{theorem}

%例题1.2.36
\begin{example}
  证明:
  \[
    \lim\limits_{n\to\infty}
    \dfrac{1^p+2^p+\cdots+n^p}{n^{p+1}}
    =\dfrac{1}{p+1}.
  \]
\end{example}
\exampleblank

\begin{solution}
  设题中 $p>0$.对
  \[
    S_n=1^p+2^p+\cdots+n^p
  \]
  与 $n^{p+1}$ 应用 Stolz 定理,得
  \[
    \lim\limits_{n\to\infty}\dfrac{S_n}{n^{p+1}}
    =\lim\limits_{n\to\infty}
    \dfrac{(n+1)^p}{(n+1)^{p+1}-n^{p+1}}.
  \]
  由 Lagrange 中值定理,存在 $\xi_n\in(n,n+1)$,使
  \[
    (n+1)^{p+1}-n^{p+1}=(p+1)\xi_n^p.
  \]
  故上式等于
  \[
    \lim\limits_{n\to\infty}\dfrac{1}{p+1}\left(\dfrac{n+1}{\xi_n}\right)^p
    =\dfrac{1}{p+1}.
  \]
\end{solution}

%例题1.2.37
\begin{example}
  计算下列极限:
  \begin{enumerate}
    \item $\displaystyle
            I=\lim\limits_{n\to\infty}\dfrac{1}{\sqrt n}\sum\limits_{k=1}^n\dfrac{a_k}{\sqrt k}
            \quad\text{(其中 }a_k\to a\ (k\to\infty)\text{)}$;
    \item $\displaystyle
            I=\lim\limits_{n\to\infty}\dfrac{\sum\limits_{k=1}^n\dfrac{1}{k}}{\ln n}$;
    \item $\displaystyle
            I=\lim\limits_{n\to\infty}\dfrac{1+a+2a^2+\cdots+na^n}{na^{n+2}}
            \quad(a>1)$.
  \end{enumerate}
\end{example}
\exampleblank

\begin{solution}
  (1) 令 $S_n=\displaystyle \sum\limits_{k=1}^n \dfrac{a_k}{\sqrt k}$.因 $\sqrt n\to+\infty$,由 Stolz 定理,
  \[
    \lim\limits_{n\to\infty}\dfrac{S_n}{\sqrt n}
    =\lim\limits_{n\to\infty}
    \dfrac{\dfrac{a_n}{\sqrt n}}{\sqrt n-\sqrt{n-1}}
    =\lim\limits_{n\to\infty}a_n\dfrac{\sqrt n+\sqrt{n-1}}{\sqrt n}=2a.
  \]

  (2) 记 $H_n=1+\dfrac{1}{2}+\cdots+\dfrac{1}{n}$.由 Stolz 定理,
  \[
    \lim\limits_{n\to\infty}\dfrac{H_n}{\ln n}
    =\lim\limits_{n\to\infty}
    \dfrac{\dfrac{1}{n+1}}{\ln(n+1)-\ln n}
    =\lim\limits_{n\to\infty}\dfrac{1}{(n+1)\ln\left(1+\dfrac{1}{n}\right)}=1.
  \]

  (3) 令 $S_n=1+a+2a^2+\cdots+na^n$.当 $a>1$ 时
  $na^{n+2}\to+\infty$,且
  \[
    \dfrac{S_{n+1}-S_n}{(n+1)a^{n+3}-na^{n+2}}
    =\dfrac{(n+1)a^{n+1}}{(n+1)a^{n+3}-na^{n+2}}
    =\dfrac{n+1}{a[(a-1)n+a]}
    \longrightarrow\dfrac{1}{a(a-1)}.
  \]
  由 Stolz 定理即得结论.
\end{solution}

%例题1.2.38
\begin{example}
  设
  \[
    S_n=\dfrac{\sum\limits_{k=0}^n\ln C_{n}^{k}}{n^2},
  \]
  求 $\displaystyle \lim\limits_{n\to\infty}S_n$.
\end{example}
\exampleblank

\begin{solution}
  为应用 Stolz 定理,记分子为
  \[
    B_n=\sum\limits_{k=0}^n\ln C_n^k.
  \]
  由
  \[
    B_n=(n+1)\ln(n!)-2\sum\limits_{k=0}^n\ln(k!),
  \]
  可直接算得
  \[
    B_{n+1}-B_n=n\ln(n+1)-\ln(n!).
  \]
  因 $n^2$ 严格递增且趋于 $+\infty$,由 Stolz 定理及 Stirling 公式
  $\ln(n!)=n\ln n-n+O(\ln n)$,
  \[
    \begin{aligned}
    \lim\limits_{n\to\infty}\dfrac{B_n}{n^2}
    &=\lim\limits_{n\to\infty}
      \dfrac{n\ln(n+1)-\ln(n!)}{(n+1)^2-n^2}\\
    &=\lim\limits_{n\to\infty}
      \dfrac{n\ln\left(1+\dfrac{1}{n}\right)+n+O(\ln n)}{2n+1}
    =\dfrac{1}{2}.
    \end{aligned}
  \]
\end{solution}

%例题1.2.39
\begin{example}
  设
  \[
    \lim\limits_{n\to\infty}\sum\limits_{k=1}^n\dfrac{1}{k^2}=m,
  \]
  求极限
  \[
    \lim\limits_{n\to\infty}
    n^2\left(m-\sum\limits_{k=1}^n\dfrac{1}{k^2}-\dfrac{1}{n}\right).
  \]
\end{example}
\exampleblank

\begin{solution}
  由 $m=\displaystyle\sum\limits_{k=1}^{\infty}\dfrac{1}{k^2}$ 及积分判别法,
  题中括号内的数列趋于 $0$.因此可与 $\dfrac{1}{n^2}$ 使用 Stolz 定理的零比零型:
  \[
    \begin{aligned}
    \lim\limits_{n\to\infty}
    n^2\left(m-\sum\limits_{k=1}^n\dfrac{1}{k^2}-\dfrac{1}{n}\right)
    =\lim\limits_{n\to\infty}
    \dfrac{
      \dfrac{1}{(n+1)^2}-\dfrac{1}{n}+\dfrac{1}{n+1}
    }{
      \dfrac{1}{n^2}-\dfrac{1}{(n+1)^2}
    }
    =\lim\limits_{n\to\infty}\left(-\dfrac{n}{2n+1}\right)
    =-\dfrac{1}{2}.
    \end{aligned}
  \]
\end{solution}

%例题1.2.40
\begin{example}
  设
  \[
    \lim\limits_{n\to\infty}n(A_n-A_{n-1})=0,
  \]
  试证:当
  \[
    \lim\limits_{n\to\infty}\dfrac{A_1+A_2+\cdots+A_n}{n}.
  \]
  存在时,
  \[
    \lim\limits_{n\to\infty}A_n
    =
    \lim\limits_{n\to\infty}\dfrac{A_1+A_2+\cdots+A_n}{n}.
  \]
\end{example}
\exampleblank

\begin{solution}
  直接作恒等变形:
  \[
    A_n-\dfrac{A_1+A_2+\cdots+A_n}{n}
    =\dfrac{1}{n}\sum\limits_{k=1}^{n-1}k(A_{k+1}-A_k).
  \]
  题设给出 $(k+1)(A_{k+1}-A_k)\to0$,故
  $k(A_{k+1}-A_k)\to0$.由 Ces\`aro 平均定理,上式右端趋于 $0$.
  因此 $A_n$ 与题中已知存在的算术平均极限相同.
\end{solution}

%例题1.2.41
\begin{example}
  设 $a_1>1$,
  \[
    a_{n+1}=a_n+\dfrac{1}{a_n},
  \]
  证明:
  \[
    \lim\limits_{n\to\infty}\dfrac{a_n}{\sqrt{2n}}=1.
  \]
\end{example}
\exampleblank

\begin{solution}
  由递推式 $a_{n+1}>a_n>0$,故 $\{a_n\}$ 单调递增.若它有有限极限 $L>0$,
  则令 $n\to\infty$ 会得到 $L=L+\dfrac{1}{L}$,矛盾.所以 $a_n\to+\infty$.

  平方递推式得
  \[
    a_{n+1}^2-a_n^2=2+\dfrac{1}{a_n^2}\longrightarrow2.
  \]
  对 $a_n^2$ 与 $n$ 使用 Stolz 定理,
  \[
    \lim\limits_{n\to\infty}\dfrac{a_n^2}{n}=2.
  \]
  由于 $a_n>0$,故
  \[
    \lim\limits_{n\to\infty}\dfrac{a_n}{\sqrt{2n}}=1.
  \]
\end{solution}

%例题1.2.42
\begin{example}
  设数列 $\{a_n\}$ 满足 $0<a_1<\dfrac{\pi}{2}$ 且
  \[
    a_{n+1}=\sin a_n,
  \]
  证明:
  \[
    \lim\limits_{n\to\infty}\sqrt{\dfrac{n}{3}}\,a_n=1.
  \]
\end{example}
\exampleblank

\begin{solution}
  在 $\left(0,\dfrac{\pi}{2}\right)$ 上有 $0<\sin x<x$,故 $0<a_{n+1}<a_n$.于是
  $a_n\to L\geq0$,而 $L=\sin L$ 迫使 $L=0$,从而 $\dfrac{1}{a_n^2}$ 严格递增且
  趋于 $+\infty$.

  要证的结论等价于 $na_n^2\to3$,而
  \[
    na_n^2=\dfrac{n}{\dfrac{1}{a_n^2}}.
  \]
  由展开式
  \[
    \sin x=x-\dfrac{x^3}{6}+O(x^5)
  \]
  可得
  \[
    \dfrac{1}{\sin^2x}-\dfrac{1}{x^2}=\dfrac{1}{3}+O(x^2)
    \qquad(x\to0).
  \]
  以 $x=a_n$ 代入,得
  \[
    \dfrac{1}{a_{n+1}^2}-\dfrac{1}{a_n^2}
    =\dfrac{1}{\sin^2a_n}-\dfrac{1}{a_n^2}
    =\dfrac{1}{3}+O(a_n^2)\longrightarrow\dfrac{1}{3}.
  \]
  对上面写出的商使用 Stolz 定理,得到
  \[
    \lim\limits_{n\to\infty}na_n^2
    =\lim\limits_{n\to\infty}
    \dfrac{1}{\dfrac{1}{a_{n+1}^2}-\dfrac{1}{a_n^2}}
    =3.
  \]
  所以
  \[
    \sqrt{\dfrac{n}{3}}\,a_n
    =\sqrt{\dfrac{na_n^2}{3}}\longrightarrow1.
  \]
\end{solution}

%例题1.2.43
\begin{example}
  设数列 $\{a_n\}$ 满足 $0<a_n<1$,
  \[
    a_{n+1}=a_n(1-a_n),
  \]
  证明:
  \begin{enumerate}
    \item $\displaystyle \lim\limits_{n\to\infty}na_n=1$;
    \item $\displaystyle \lim\limits_{n\to\infty}\dfrac{n(1-na_n)}{\ln n}=1$.
  \end{enumerate}
\end{example}
\exampleblank

\begin{solution}
  由 $0<a_n<1$ 可知 $0<a_{n+1}<a_n$,故 $a_n\to L\geq0$.
  在递推式中令 $n\to\infty$,得 $L=L(1-L)$,所以 $L=0$,从而 $\dfrac{1}{a_n}$
  严格递增且趋于 $+\infty$.

  (1) 把 $na_n$ 写成商
  \[
    na_n=\dfrac{n}{\dfrac{1}{a_n}},
  \]
  再对该商使用 Stolz 定理.由
  \[
    \dfrac{1}{a_{n+1}}-\dfrac{1}{a_n}
    =\dfrac{1}{a_n(1-a_n)}-\dfrac{1}{a_n}
    =\dfrac{1}{1-a_n}\longrightarrow1,
  \]
  得
  \[
    \lim\limits_{n\to\infty}na_n
    =\lim\limits_{n\to\infty}\dfrac{n}{\dfrac{1}{a_n}}
    =\lim\limits_{n\to\infty}
    \dfrac{1}{\dfrac{1}{a_{n+1}}-\dfrac{1}{a_n}}
    =\lim\limits_{n\to\infty}(1-a_n)=1.
  \]

  (2) 把所求表达式写成两个因子之积
  \[
    \dfrac{n(1-na_n)}{\ln n}
    =\dfrac{n}{\dfrac{1}{a_n}}\cdot
    \dfrac{\dfrac{1}{a_n}-n}{\ln n},
  \]
  第一个因子 $\dfrac{n}{\dfrac{1}{a_n}}=na_n$ 由 (1) 趋于 $1$,第二个因子用
  Stolz 定理,分母$\ln n$单调递增趋于 $+\infty$.由
  \[
    \left(\dfrac{1}{a_{n+1}}-(n+1)\right)-\left(\dfrac{1}{a_n}-n\right)
    =\dfrac{1}{1-a_n}-1=\dfrac{a_n}{1-a_n},
  \]
  得
  \[
    \lim\limits_{n\to\infty}\dfrac{\dfrac{1}{a_n}-n}{\ln n}
    =\lim\limits_{n\to\infty}
    \dfrac{\dfrac{a_n}{1-a_n}}{\ln\left(1+\dfrac{1}{n}\right)}
    =\lim\limits_{n\to\infty}
    \dfrac{na_n}{1-a_n}\cdot\dfrac{1}{n\ln\left(1+\dfrac{1}{n}\right)}
    =1,
  \]
  因此
  \[
    \lim\limits_{n\to\infty}\dfrac{n(1-na_n)}{\ln n}=1\cdot1=1.
  \]
\end{solution}

%例题1.2.44
\begin{example}
  设 $0<a_1<1$,
  \[
    a_{n+1}=\ln(1+a_n),
  \]
  则
  \[
    \lim\limits_{n\to\infty}\dfrac{n(na_n-2)}{\ln n}=\dfrac{2}{3}.
  \]
\end{example}
\exampleblank

\begin{solution}
  因为 $0<\ln(1+x)<x$ 对 $x>0$ 成立,所以
  $0<a_{n+1}<a_n$,从而 $a_n\to L\geq0$.由 $L=\ln(1+L)$ 得 $L=0$,
  于是 $\dfrac{1}{a_n}$ 严格递增且趋于 $+\infty$.

  先把 $na_n$ 写成商
  \[
    na_n=\dfrac{n}{\dfrac{1}{a_n}},
  \]
  再用 Stolz 定理.当 $x\to0$ 时,
  \[
    \dfrac{1}{\ln(1+x)}
    =\dfrac{1}{x}+\dfrac{1}{2}-\dfrac{x}{12}+O(x^2).
  \]
  以 $x=a_n$ 代入,得
  \[
    \dfrac{1}{a_{n+1}}-\dfrac{1}{a_n}
    =\dfrac{1}{\ln(1+a_n)}-\dfrac{1}{a_n}
    =\dfrac{1}{2}-\dfrac{a_n}{12}+O(a_n^2)\longrightarrow\dfrac{1}{2},
  \]
  故
  \[
    \lim\limits_{n\to\infty}na_n
    =\lim\limits_{n\to\infty}\dfrac{n}{\dfrac{1}{a_n}}
    =\lim\limits_{n\to\infty}
    \dfrac{1}{\dfrac{1}{a_{n+1}}-\dfrac{1}{a_n}}
    =2.
  \]

  再把所求表达式写成两个因子之积
  \[
    \dfrac{n(na_n-2)}{\ln n}
    =\dfrac{n}{\dfrac{1}{a_n}}\cdot
    \dfrac{n-\dfrac{2}{a_n}}{\ln n},
  \]
  第一个因子由上面的结果趋于 $2$,第二个因子用 Stolz 定理,分母 $\ln n$
  单调递增趋于 $+\infty$.由上面的展开式,
  \[
    \left(n+1-\dfrac{2}{a_{n+1}}\right)-\left(n-\dfrac{2}{a_n}\right)
    =1-2\left(\dfrac{1}{\ln(1+a_n)}-\dfrac{1}{a_n}\right)
    =\dfrac{a_n}{6}+O(a_n^2),
  \]
  于是
  \[
    \lim\limits_{n\to\infty}\dfrac{n-\dfrac{2}{a_n}}{\ln n}
    =\lim\limits_{n\to\infty}
    \dfrac{\dfrac{a_n}{6}+O(a_n^2)}{\ln\left(1+\dfrac{1}{n}\right)}
    =\lim\limits_{n\to\infty}
    \dfrac{\dfrac{na_n}{6}+O(na_n^2)}
    {n\ln\left(1+\dfrac{1}{n}\right)}
    =\dfrac{1}{3},
  \]
  其中最后一步用了 $na_n\to2$.因此
  \[
    \lim\limits_{n\to\infty}\dfrac{n(na_n-2)}{\ln n}
    =2\cdot\dfrac{1}{3}=\dfrac{2}{3}.
  \]
\end{solution}



\methodsection{(E) 利用函数极限}

结合海涅归结原理,可将数列极限问题转化为函数极限问题,而函数极限求法留在下章重点介绍.


\newpage

\section{递推数列}

假设用某种方法证明了递推数列的极限存在,则在递推公式里取极限,便可得极限值 $A$ 应满足的方程.当然,若能求得数列通项,极限值当然可以直接求得,求通项技巧高中大多已经学习过,就不再介绍了,接下来介绍几种不需要通项就可得到极限存在性的常用方法.



\subsection{单调有界原理}

\noindent\textbf{(\markstar)} 单调有界原理是最常见也最常用的方法,要分别证明单调性与有界性.单调性通常有以下几种判断方法:

\begin{enumerate}
  \item 对任意 $n\in\mathbb{N}$,若 $x_{n+1}-x_n\geq0$,则 $x_n$ 单增;若 $x_{n+1}-x_n\leq0$,则 $x_n$ 单减;
  \item 对任意 $n\in\mathbb{N}$,若 $\dfrac{x_{n+1}}{x_n}\geq1$,则 $x_n$ 单增;若 $\dfrac{x_{n+1}}{x_n}\leq1$,则 $x_n$ 单减;
  \item 若 $x_{n+1}=f(x_n)$,$f'(x)\geq0$,则若 $x_1\leq x_2$,则 $x_n$ 单增;若 $x_1\geq x_2$,则 $x_n$ 单减.
\end{enumerate}

此外,数学归纳法及常用不等式都是证明单调,有界的重要工具.

%例题1.3.1
\begin{example}
  证明数列
  \[
    x_0=1,\qquad x_{n+1}=\sqrt{2x_n},\qquad n=0,1,2,\ldots.
  \]
  有极限,并求其值.
\end{example}
\exampleblank

\begin{solution}
  由 $x_0=1<2$.若 $0<x_n<2$,则
  \[
    x_n<x_{n+1}=\sqrt{2x_n}<2,
  \]
  其中第一式等价于 $x_n<2$.故归纳可知 $\{x_n\}$ 单调递增且以 $2$ 为上界.
  设 $x_n\to L$,则 $L\geq1$,并且
  \[
    L=\sqrt{2L}.
  \]
  所以 $L=2$.即 $x_n\to2$.
\end{solution}

%例题1.3.2
\begin{example}[(\markstar)]
  已知连续函数 $f(x)$ 在 $\lbrack 1,+\infty)$ 单调递减,且 $f(x)>0$,证明数列
  \[
    a_n=\sum\limits_{k=1}^n f(k)-\int_1^n f(x)\,\mathrm{d}x.
  \]
  收敛.
\end{example}
\exampleblank

\begin{solution}
  由于 $f$ 单调递减,
  \[
    a_{n+1}-a_n
    =f(n+1)-\int_n^{n+1}f(x)\,\mathrm{d}x\leq0.
  \]
  又
  \[
    \int_1^n f(x)\,\mathrm{d}x
    \leq\sum\limits_{k=1}^{n-1}f(k),
  \]
  所以 $a_n\geq f(n)>0$.因此 $\{a_n\}$ 单调递减且有下界,从而极限存在.
\end{solution}

\begin{remark}
  \begin{enumerate}
    \item 取 $f(x)=\dfrac{1}{x}$,则有
          \[
            a_n=1+\dfrac{1}{2}+\cdots+\dfrac{1}{n}-\ln n
          \]
          是收敛的;
    \item 取 $f(x)=\sqrt{\dfrac{1}{x}}$,则有
          \[
            a_n=1+\sqrt{\dfrac{1}{2}}+\cdots+\sqrt{\dfrac{1}{n}}-2\sqrt n
          \]
          是收敛的;
    \item 取 $f(x)=\dfrac{1}{x\ln x}$,则有
          \[
            a_n=\dfrac{1}{2\ln2}+\dfrac{1}{3\ln3}+\cdots+\dfrac{1}{n\ln n}-\ln\ln n
          \]
          是收敛的.
  \end{enumerate}
\end{remark}

%例题1.3.3
\begin{example}
  设 $A>0$,$a_1>0$,
  \[
    a_{n+1}=\dfrac{1}{2}\left(a_n+\dfrac{A}{a_n}\right),
  \]
  证明:$\displaystyle \lim\limits_{n\to\infty}a_n$ 存在,并求其值.
\end{example}
\exampleblank

\begin{solution}
  由算术--几何平均值不等式,
  \[
    a_2=\dfrac{1}{2}\left(a_1+\dfrac{A}{a_1}\right)\geq\sqrt A.
  \]
  若 $a_n\geq\sqrt A$,则
  \[
    a_{n+1}-\sqrt A=\dfrac{(a_n-\sqrt A)^2}{2a_n}\geq0,
  \]
  且
  \[
    a_{n+1}-a_n=\dfrac{A-a_n^2}{2a_n}\leq0.
  \]
  故 $\{a_n\}_{n\geq2}$ 单调递减且下界为 $\sqrt A$,于是存在极限 $L\geq\sqrt A$.
  由
  \[
    L=\dfrac{1}{2}\left(L+\dfrac{A}{L}\right)
  \]
  得 $L^2=A$.因此 $a_n\to\sqrt A$.
\end{solution}

%例题1.3.4
\begin{example}
  设 $c>0$,$0<x_1<\dfrac{1}{c}$,且
  \[
    x_{n+1}=x_n(2-cx_n),
  \]
  证明:$\{x_n\}$ 收敛并求其极限.
\end{example}
\exampleblank

\begin{solution}
  先证 $0<x_n<\dfrac{1}{c}$.若该式对某个 $n$ 成立,则
  \[
    x_{n+1}=x_n(2-cx_n)>0,
  \]
  并且
  \[
    \dfrac{1}{c}-x_{n+1}
    =c\left(\dfrac{1}{c}-x_n\right)^2>0.
  \]
  故归纳可知该区间对所有 $n$ 不变.又
  \[
    x_{n+1}-x_n
    =cx_n\left(\dfrac{1}{c}-x_n\right)>0.
  \]
  所以 $\{x_n\}$ 单调递增且以 $\dfrac{1}{c}$ 为上界.设其极限为 $L>0$,则
  $L=L(2-cL)$,从而 $L=\dfrac{1}{c}$.
\end{solution}

%例题1.3.5
\begin{example}
  已知数列 $\{a_n\},\{b_n\}$ 满足:$a_1>b_1>0$,
  \[
    a_{n+1}=\dfrac{a_n+b_n}{2},
    \qquad
    b_{n+1}=\sqrt{a_nb_n},
  \]
  证明:$\displaystyle \lim\limits_{n\to\infty}a_n$ 与 $\displaystyle \lim\limits_{n\to\infty}b_n$ 都存在且相等.
\end{example}
\exampleblank

\begin{solution}
  由算术--几何平均值不等式,
  \[
    b_n\leq a_n,
    \qquad
    b_n\leq b_{n+1}=\sqrt{a_nb_n}\leq a_{n+1}=\dfrac{a_n+b_n}{2}\leq a_n.
  \]
  所以 $\{a_n\}$ 单调递减,$\{b_n\}$ 单调递增,且始终有 $b_n\leq a_n$.
  因此存在 $\alpha,\beta>0$,使 $a_n\to\alpha$,$b_n\to\beta$,并且
  $\beta\leq\alpha$.在
  \[
    a_{n+1}=\dfrac{a_n+b_n}{2}
  \]
  中令 $n\to\infty$,得 $\alpha=\dfrac{\alpha+\beta}{2}$,故 $\alpha=\beta$.
  所以两数列收敛于同一正数.
\end{solution}

%例题1.3.6
\begin{example}
  已知数列 $\{a_n\}$ 满足 $0<a_n<2$ 且
  \[
    (2-a_n)a_{n+1}\geq1,
  \]
  证明:$\{a_n\}$ 收敛,并求其极限.
\end{example}
\exampleblank

\begin{solution}
  由题设及 $0<a_n<2$,
  \[
    a_{n+1}\geq\dfrac{1}{2-a_n}.
  \]
  又
  \[
    \dfrac{1}{2-a_n}-a_n
    =\dfrac{(a_n-1)^2}{2-a_n}\geq0,
  \]
  故 $a_{n+1}\geq a_n$.因此 $\{a_n\}$ 单调递增且以 $2$ 为上界,存在极限
  $L\leq2$.若 $L=2$,则 $(2-a_n)a_{n+1}\to0$,与题设矛盾,故 $L<2$.
  令 $n\to\infty$,得到
  \[
    (2-L)L\geq1,
  \]
  即 $-(L-1)^2\geq0$.所以 $L=1$,即 $a_n\to1$.
\end{solution}

%例题1.3.7
\begin{example}
  设数列 $\{x_n\}$ 满足 $x_n>0$ 且
  \[
    x_{n+1}+\dfrac{4}{x_n}<4,
  \]
  证明:$\{x_n\}$ 收敛并求其极限.
\end{example}
\exampleblank

\begin{solution}
  因 $x_{n+1}>0$,题设给出 $\dfrac{4}{x_n}<4$,故 $x_n>1$.又
  \[
    4-\dfrac{4}{x_n}-x_n
    =-\dfrac{(x_n-2)^2}{x_n}\leq0,
  \]
  所以
  \[
    x_{n+1}<4-\dfrac{4}{x_n}\leq x_n.
  \]
  于是 $\{x_n\}$ 单调递减且下界为 $1$,故 $x_n\to L\geq1$.
  对原不等式取极限上界,得
  \[
    L+\dfrac{4}{L}\leq4.
  \]
  这等价于 $(L-2)^2\leq0$,故 $L=2$.
\end{solution}



\subsection{压缩映像原理}

\begin{theorem}[\markstar]
  对于任一数列 $\{x_n\}$ 而言,若存在常数 $r$,使得对任意 $n\in\mathbb{N}$,恒有
  \[
    |x_{n+1}-x_n|\leq r|x_n-x_{n-1}|,
    \qquad 0<r<1,
  \]
  则数列 $\{x_n\}$ 收敛.
\end{theorem}

\begin{corollary}[\markstar]

    设 $I\subseteq\mathbb{R}$ 为非空闭集,$f:I\to I$.

    若存在常数 $0<r<1$,使得
    \[
        |f(x)-f(y)|\leq r|x-y|,
        \qquad x,y\in I,
    \]
    则 $f$ 在 $I$ 上存在唯一不动点 $x^*$.

    若进一步 $f$ 在 $I$ 的内部可微且
    \[
        |f'(x)|\leq r,
    \]
    则上述条件成立.

\end{corollary}

%例题1.3.8
\begin{example}
  证明:
  \begin{enumerate}
    \item 存在唯一的 $c\in(0,1)$,使得 $c=e^{-c}$;
    \item 设 $x_1\in(0,1)$,$x_{n+1}=e^{-x_n}$,则
          $\displaystyle \lim\limits_{n\to\infty}x_n=c$.
  \end{enumerate}
\end{example}
\exampleblank

\begin{solution}
  (1) 令 $h(x)=x-e^{-x}$.则
  \[
    h(0)=-1<0,
    \qquad h(1)=1-e^{-1}>0,
    \qquad h'(x)=1+e^{-x}>0.
  \]
  由介值定理和严格单调性,存在唯一的 $c\in(0,1)$ 使 $c=e^{-c}$.

  (2) 令 $f(x)=e^{-x}$,$I=\lbrack e^{-1},e^{-e^{-1}}\rbrack$.由 $x_1\in(0,1)$ 可知
  $x_n\in(0,1)$;进一步 $x_2\in(e^{-1},1)$,故 $x_3\in I$.
  又 $f$ 在 $I$ 上递减,而
  \[
    f(e^{-1})=e^{-e^{-1}},
    \qquad f(e^{-e^{-1}})=e^{-e^{-e^{-1}}}>e^{-1},
  \]
  故 $f(I)\subseteq I$.由于 $f$ 在 $I$ 上连续,在 $I$ 的内部可微,且对一切内点 $x$ 有
  \[
    |f'(x)|=e^{-x}\leq e^{-e^{-1}}=:q<1,
  \]
  由推论 1.3.1,$\{x_n\}_{n\geq3}$ 收敛,从而 $\{x_n\}$ 收敛,且收敛于 $f$ 在 $I$ 上的
  唯一不动点 $x^*$.因 $I\subseteq(0,1)$,由 (1) 中 $(0,1)$ 内不动点的唯一性得
  $x^*=c$.故 $x_n\to c$.
\end{solution}

%例题1.3.9
\begin{example}
  设 $x_0>0$,
  \[
    x_n=3+\dfrac{1}{\sqrt{x_{n-1}}}\qquad(n=1,2,\ldots),
  \]
  证明:$\{x_n\}$ 收敛并求其极限.
\end{example}
\exampleblank

\begin{solution}
  令 $f(x)=3+x^{-\frac{1}{2}}$,$I=\lbrack3,+\infty)$.由 $x_0>0$ 得
  $x_1=3+\dfrac{1}{\sqrt{x_0}}>3$,即 $x_1\in I$;而当 $x\geq3$ 时
  \[
    f(x)=3+\dfrac{1}{\sqrt x}\in(3,3+3^{-\frac{1}{2}}\rbrack\subseteq I,
  \]
  即 $f(I)\subseteq I$.由于 $f$ 在 $I$ 上连续,在 $I$ 的内部可微,且对一切内点 $x$ 有
  \[
    0<f'(x)=\dfrac{1}{2x^{\frac{3}{2}}}\leq\dfrac{1}{2\cdot3^{\frac{3}{2}}}=:q<1,
  \]
  由推论 1.3.1,$\{x_n\}$ 收敛,且收敛于 $f$ 在 $I$ 上的唯一不动点 $x^*$.令
  $t=\sqrt{x^*}$,由
  \[
    x^*=3+\dfrac{1}{\sqrt{x^*}}
  \]
  得 $t^3-3t-1=0$.因 $x^*\geq3$,有 $t\geq\sqrt3>0$,而该方程的唯一正根为
  $t=2\cos\left(\dfrac{\pi}{9}\right)$,所以
  \[
    x^*=4\cos^2\dfrac{\pi}{9}.
  \]
\end{solution}

%例题1.3.10
\begin{example}
  设
  \[
    u_1=3,
    \qquad
    u_2=3+\dfrac{4}{3},
    \qquad
    u_3=3+\dfrac{4}{3+\dfrac{4}{3}},
    \qquad\ldots.
  \]
  如果数列 $\{u_n\}$ 收敛,计算其极限,并证明 $\{u_n\}$ 收敛于上述极限.
\end{example}
\exampleblank

\begin{solution}
  该连分数满足
  \[
    u_{n+1}=3+\dfrac{4}{u_n},
    \qquad u_1=3.
  \]
  令 $f(x)=3+\dfrac{4}{x}$,$I=\lbrack3,+\infty)$.由 $u_1=3\in I$,而当 $x\geq3$ 时
  \[
    f(x)=3+\dfrac{4}{x}\in\left(3,\dfrac{13}{3}\right\rbrack\subseteq I,
  \]
  即 $f(I)\subseteq I$.由于 $f$ 在 $I$ 上连续,在 $I$ 的内部可微,且对一切内点 $x$ 有
  \[
    |f'(x)|=\dfrac{4}{x^2}\leq\dfrac{4}{9}<1,
  \]
  由推论 1.3.1,$\{u_n\}$ 收敛,且收敛于 $f$ 在 $I$ 上的唯一不动点 $x^*$.由
  \[
    x^*=3+\dfrac{4}{x^*}
  \]
  得 $(x^*)^2-3x^*-4=0$,其正根为 $x^*=4$,故 $u_n\to4$.
\end{solution}

%例题1.3.11
\begin{example}
  设 $x_1=1$,
  \[
    x_{n+1}=\dfrac{1+2x_n}{1+x_n}\qquad(n=1,2,\ldots),
  \]
  证明:$\{x_n\}$ 收敛,并求 $\displaystyle \lim\limits_{n\to\infty}x_n$.
\end{example}
\exampleblank

\begin{solution}
  令 $f(x)=\dfrac{1+2x}{1+x}$,$I=\lbrack1,+\infty)$.由 $x_1=1\in I$,而当 $x\geq1$ 时
  \[
    f(x)=2-\dfrac{1}{1+x}\in\left\lbrack\dfrac{3}{2},2\right)\subseteq I,
  \]
  即 $f(I)\subseteq I$.由于 $f$ 在 $I$ 上连续,在 $I$ 的内部可微,且对一切内点 $x$ 有
  \[
    0<f'(x)=\dfrac{1}{(1+x)^2}\leq\dfrac{1}{4}<1,
  \]
  由推论 1.3.1,$\{x_n\}$ 收敛,且收敛于 $f$ 在 $I$ 上的唯一不动点 $x^*$.由
  \[
    x^*=\dfrac{1+2x^*}{1+x^*}
  \]
  得 $(x^*)^2-x^*-1=0$,故 $x^*=\dfrac{1+\sqrt5}{2}$,即 $x_n\to\dfrac{1+\sqrt5}{2}$.
\end{solution}



\subsection{不动点法}

事实上,前面两种方法已知能解决绝大部分问题,但对于部分难题,不动点法尤为有效.

\begin{proposition}[\markstar]
  已知数列 $\{x_n\}$ 在区间 $I$ 上由
  \[
    x_{n+1}=f(x_n)\qquad(n=1,2,\ldots).
  \]
  给出,$f$ 是 $I$ 上连续增函数.若 $f$ 在 $I$ 上有不动点 $x^*$ 满足
  \[
    (x_1-f(x_1))(x_1-x^*)\geq0,
  \]
  则此时数列 $\{x_n\}$ 必收敛,且极限 $A$ 满足 $A=f(A)$.若上式中的``$\geq$''改为``$>$'',则 $x^*$ 为唯一不动点.

  特别地,若 $f$ 可导且 $0<f'(x)<1$,则 $f$ 严格递增,且上式自然成立,这时 $f$ 的不动点存在且唯一,从而 $A=x^*$.
\end{proposition}

\begin{proposition}[\markstar]
  设数列 $\{x_n\}$ 由 $x_{n+1}=f(x_n)$ 给出,递推下去取范围不超过 $I$,$f(x)$ 在 $I$ 上有 $f'(x)<0$($f$ 严格递减).若 $f$ 有不动点 $x^*$,则 $\{x_{2n}\},\{x_{2n\pm1}\}$ 分别在 $x^*$ 之两侧,且 $\{x_{2n}\},\{x_{2n-1}\}$ 均单调,并且具有相反单调性.
\end{proposition}

\begin{remark}
  这时 $F(x)=f(f(x))$ 是它们的递推函数,且
  \[
    F'(x)=f'(f(x))f'(x)>0.
  \]
  因此,可按命题 1.3.1 分别处理 $\{x_{2n}\}$ 与 $\{x_{2n-1}\}$.
\end{remark}

%例题1.3.12
\begin{example}
  设 $x_1>0$,
  \[
    x_{n+1}=\dfrac{c(1+x_n)}{c+x_n}\qquad(c>1),
  \]
  求 $\displaystyle \lim\limits_{n\to\infty}x_n$.
\end{example}
\exampleblank

\begin{solution}

  令
  \[
    f(x)=\dfrac{c(1+x)}{c+x}.
  \]
  因 $x_1>0$,由递推式知 $x_n>0$.正不动点满足
  \[
    x=\dfrac{c(1+x)}{c+x},
  \]
  即 $x^2=c$,故正不动点为 $s=\sqrt c$.

  直接计算得
  \[
    x_{n+1}-s
    =\dfrac{s(s-1)}{c+x_n}(x_n-s).
  \]
  因此
  \[
    |x_{n+1}-s|
    \leq\dfrac{s-1}{s}|x_n-s|.
  \]
  由于 $0<\dfrac{s-1}{s}<1$,反复迭代得
  \[
    |x_n-s|
    \leq
    \left(\dfrac{s-1}{s}\right)^{n-1}|x_1-s|
    \longrightarrow0.
  \]
  故
  \[
    \lim\limits_{n\to\infty}x_n=\sqrt c.
  \]

\end{solution}

%例题1.3.13
\begin{example}
  设
  \[
    x_1=1,
    \qquad
    x_2=\dfrac{1}{2},
    \qquad
    x_{n+1}=\dfrac{1}{1+x_n},
  \]
  求 $\displaystyle \lim\limits_{n\to\infty}x_n$.
\end{example}
\exampleblank

\begin{solution}

  令
  \[
    f(x)=\dfrac{1}{1+x}.
  \]
  从 $x_2=\dfrac{1}{2}$ 开始,
  \[
    \dfrac{1}{2}\leq x_n\leq\dfrac{2}{3},
  \]
  且该区间在 $f$ 下不变.

  正不动点满足
  \[
    x=\dfrac{1}{1+x},
  \]
  故
  \[
    \alpha=\dfrac{\sqrt5-1}{2}.
  \]
  由 $\alpha=\dfrac{1}{1+\alpha}$,
  \[
    x_{n+1}-\alpha
    =-\dfrac{x_n-\alpha}{(1+x_n)(1+\alpha)}.
  \]
  因此
  \[
    |x_{n+1}-\alpha|
    \leq\dfrac{2}{3}|x_n-\alpha|.
  \]
  反复迭代即得
  \[
    x_n\longrightarrow\alpha
    =\dfrac{\sqrt5-1}{2}.
  \]

\end{solution}

%例题1.3.14
\begin{example}
  设 $f(x)$ 映 $\lbrack a,b\rbrack$ 为自身,且
  \[
    |f(x)-f(y)|\leq|x-y|.
  \]
  任取 $x_1\in\lbrack a,b\rbrack$,令
  \[
    x_{n+1}=\dfrac{1}{2}[x_n+f(x_n)],
  \]
  求证数列有极限 $x^*$,且 $x^*$ 满足方程
  \[
    f(x^*)=x^*.
  \]
\end{example}
\exampleblank

\begin{solution}

  令
  \[
    g(x)=\dfrac{x+f(x)}{2}.
  \]
  因 $f$ 把 $\lbrack a,b\rbrack$ 映入自身,故 $g$ 也把
  $\lbrack a,b\rbrack$ 映入自身.

  对任意 $x<y$,由
  \[
    |f(y)-f(x)|\leq y-x
  \]
  可得
  \[
    f(y)-f(x)\geq-(y-x),
  \]
  从而
  \[
    g(y)-g(x)
    =\dfrac{(y-x)+(f(y)-f(x))}{2}\geq0.
  \]
  因此 $g$ 单调不减.

  若 $x_2\geq x_1$,则由 $x_{n+1}=g(x_n)$ 及 $g$ 的单调性,
  归纳可得
  \[
    x_1\leq x_2\leq x_3\leq\cdots.
  \]
  若 $x_2\leq x_1$,则同理
  \[
    x_1\geq x_2\geq x_3\geq\cdots.
  \]
  又因为 $x_n\in\lbrack a,b\rbrack$,故无论哪种情形,
  $\{x_n\}$ 均为有界单调数列,从而存在极限.设
  \[
    x_n\to x^*.
  \]

  由条件知 $f$ 连续,在递推式
  \[
    x_{n+1}=\dfrac{x_n+f(x_n)}{2}
  \]
  中令 $n\to\infty$,得到
  \[
    x^*=\dfrac{x^*+f(x^*)}{2},
  \]
  因而
  \[
    f(x^*)=x^*.
  \]

\end{solution}

%例题1.3.15
\begin{example}
  证明:数列 $x_0>0$,
  \[
    x_{n+1}=\dfrac{x_n(x_n^2+3a)}{3x_n^2+a}\qquad(a\geq0).
  \]
  的极限存在,并求其极限.
\end{example}
\exampleblank

\begin{solution}
  若 $a=0$,则递推式化为 $x_{n+1}=\dfrac{1}{3}x_n$,故 $x_n\to0$.

  以下设 $a>0$.正不动点满足
  \[
    x=\dfrac{x(x^2+3a)}{3x^2+a},
  \]
  因 $x>0$,约去 $x$ 后得到 $3x^2+a=x^2+3a$,故唯一的正不动点为
  $x=\sqrt a$.

  直接计算得
  \[
    \dfrac{x_{n+1}-\sqrt a}{x_{n+1}+\sqrt a}
    =\left(\dfrac{x_n-\sqrt a}{x_n+\sqrt a}\right)^3.
  \]
  迭代可得
  \[
    \dfrac{x_n-\sqrt a}{x_n+\sqrt a}
    =\left(\dfrac{x_0-\sqrt a}{x_0+\sqrt a}\right)^{3^n}
    \longrightarrow0,
  \]
  其中底数的绝对值小于 $1$.将上式解出 $x_n$,得到
  \[
    x_n=\sqrt a\,
    \dfrac{
      1+\left(\dfrac{x_0-\sqrt a}{x_0+\sqrt a}\right)^{3^n}
    }{
      1-\left(\dfrac{x_0-\sqrt a}{x_0+\sqrt a}\right)^{3^n}
    }
    \longrightarrow\sqrt a.
  \]
\end{solution}



\subsection{利用上下极限 *}

这部分方法由于预备知识理论性较强,故放在本章最后一节作为选学内容.


\newpage

\section{上,下极限 *}

本节通过几个等价定义由不同角度来刻画上,下极限,其中最直观的是用极限点来定义:

\begin{definition}
  数列的收敛子列的极限称为数列的极值点,而数列的最大(小)极值点称为数列的上(下)极限,记为
  $\displaystyle \varlimsup\limits_{n\to\infty}x_n$
  ($\displaystyle \varliminf\limits_{n\to\infty}x_n$).
\end{definition}

\begin{proposition}[\markstar\ 等价定义 1]
  每个数列 $\{x_n\}$ 都存在上极限和下极限,且成立公式:
  \[
    \varlimsup\limits_{n\to\infty}x_n
    =
    \lim\limits_{n\to\infty}\sup_{k\geq n}\{x_k\},
  \]
  \[
    \varliminf\limits_{n\to\infty}x_n
    =
    \lim\limits_{n\to\infty}\inf_{k\geq n}\{x_k\}.
  \]
\end{proposition}

\begin{proposition}[\markstar\ 等价定义 2]
  \leavevmode
  \begin{enumerate}
    \item 有限数 $b$ 是数列 $x_n$ 的下极限,当且仅当 $\{x_n\}$ 满足:
          \begin{enumerate}[label=(\roman*)]
            \item 对任意 $\varepsilon>0$,$(b-\varepsilon,b+\varepsilon)$ 中有数列的无穷多项;
            \item 存在 $N$,当 $n>N$ 时,$x_n>b-\varepsilon$;
          \end{enumerate}
    \item $+\infty$ 是 $\{x_n\}$ 的下极限,当且仅当 $\{x_n\}$ 是正无穷大量;
    \item $-\infty$ 是 $\{x_n\}$ 的下极限,当且仅当 $\{x_n\}$ 无下界.
  \end{enumerate}
  上极限同理.
\end{proposition}

\begin{theorem}[\markstar]
  上,下极限满足如下性质:
  \begin{enumerate}
    \item
          \[
            \varlimsup\limits_{n\to\infty}(-a_n)=-\varliminf\limits_{n\to\infty}a_n,
            \qquad
            \varliminf\limits_{n\to\infty}(-a_n)=-\varlimsup\limits_{n\to\infty}a_n;
          \]
    \item 若 $a_n\leq b_n$,则
          \[
            \varlimsup\limits_{n\to\infty}a_n\leq\varlimsup\limits_{n\to\infty}b_n,
            \qquad
            \varliminf\limits_{n\to\infty}a_n\leq\varliminf\limits_{n\to\infty}b_n;
          \]
    \item
          \[
            \varlimsup\limits_{n\to\infty}(a_n+b_n)
            \leq
            \varlimsup\limits_{n\to\infty}a_n+
            \varlimsup\limits_{n\to\infty}b_n,
          \]
          \[
            \varliminf\limits_{n\to\infty}(a_n+b_n)
            \geq
            \varliminf\limits_{n\to\infty}a_n+
            \varliminf\limits_{n\to\infty}b_n;
          \]
    \item
          \[
            \varliminf\limits_{n\to\infty}(a_n+b_n)
            \leq
            \varliminf\limits_{n\to\infty}a_n+
            \varlimsup\limits_{n\to\infty}b_n
            \leq
            \varlimsup\limits_{n\to\infty}(a_n+b_n);
          \]
    \item 若 $a_n,b_n>0$,则
          \[
            \varlimsup\limits_{n\to\infty}a_nb_n
            \leq
            \left(\varlimsup\limits_{n\to\infty}a_n\right)
            \left(\varlimsup\limits_{n\to\infty}b_n\right),
          \]
          \[
            \varliminf\limits_{n\to\infty}a_nb_n
            \geq
            \left(\varliminf\limits_{n\to\infty}a_n\right)
            \left(\varliminf\limits_{n\to\infty}b_n\right);
          \]
    \item 若 $a_n>0$ 且 $\displaystyle \varliminf\limits_{n\to\infty}a_n>0$,则
          \[
            \varlimsup\limits_{n\to\infty}\dfrac{1}{a_n}
            =
            \dfrac{1}{\displaystyle \varliminf\limits_{n\to\infty}a_n};
          \]
    \item 若 $\{b_n\}$ 收敛,则
          \[
            \varliminf\limits_{n\to\infty}(a_n+b_n)
            =
            \varliminf\limits_{n\to\infty}a_n+\lim\limits_{n\to\infty}b_n,
          \]
          \[
            \varlimsup\limits_{n\to\infty}(a_n+b_n)
            =
            \varlimsup\limits_{n\to\infty}a_n+\lim\limits_{n\to\infty}b_n.
          \]
  \end{enumerate}
\end{theorem}

\begin{proposition}[\markstar]
  设 $\{a_n\}$ 为正数列,则有
  \[
    \varliminf\limits_{n\to\infty}\dfrac{a_{n+1}}{a_n}
    \leq
    \varliminf\limits_{n\to\infty}\sqrt[n]{a_n}
    \leq
    \varlimsup\limits_{n\to\infty}\sqrt[n]{a_n}
    \leq
    \varlimsup\limits_{n\to\infty}\dfrac{a_{n+1}}{a_n}.
  \]
\end{proposition}

%例题1.4.1
\begin{example}
  请求下列数列 $\{a_n\}$ 的上,下极限:
  \begin{enumerate}
    \item $\displaystyle a_n=\sqrt[n]{1+2^{n(-1)^n}}$;
    \item $\displaystyle a_n=1+\dfrac{n-\cos\dfrac{n\pi}{2}}{n+1}$;
    \item $\displaystyle a_n=\dfrac{n}{3}-\left[\dfrac{n}{3}\right]$.
  \end{enumerate}
\end{example}
\exampleblank

\begin{solution}
  (1) 当 $n$ 为偶数时,
  \[
    a_n=\sqrt[n]{1+2^n}\longrightarrow2;
  \]
  当 $n$ 为奇数时,
  \[
    a_n=\sqrt[n]{1+2^{-n}}\longrightarrow1.
  \]
  故
  \[
    \varlimsup\limits_{n\to\infty}a_n=2,
    \qquad
    \varliminf\limits_{n\to\infty}a_n=1.
  \]

  (2) 有
  \[
    a_n=2-\dfrac{1+\cos\dfrac{n\pi}{2}}{n+1}\longrightarrow2.
  \]
  故上,下极限均为 $2$.

  (3) $a_n$ 是 $\dfrac{n}{3}$ 的小数部分,依次循环取 $\dfrac{1}{3},\dfrac{2}{3},0$.故
  \[
    \varlimsup\limits_{n\to\infty}a_n=\dfrac{2}{3},
    \qquad
    \varliminf\limits_{n\to\infty}a_n=0.
  \]
\end{solution}

%例题1.4.2
\begin{example}
  设 $\{a_n\}$ 满足
  \[
    \lim\limits_{n\to\infty}(a_n-2a_{n+1})=A,
  \]
  证明:$\{a_n\}$ 收敛并求其极限.
\end{example}
\exampleblank

\begin{solution}
  因为 $a_n-2a_{n+1}$ 收敛,所以存在 $M>0$,使
  $|a_n-2a_{n+1}|\leq M$.于是
  \[
    |a_{n+1}|\leq\dfrac{1}{2}|a_n|+\dfrac{M}{2},
  \]
  迭代后得到 $|a_n|\leq|a_1|+M$,故 $\{a_n\}$ 有界.

  记其下极限和上极限分别为 $\alpha,\beta$.利用移位不改变上,下极限,以及
  $a_n-2a_{n+1}\to A$,在恒等式
  \[
    a_{n+1}=\dfrac{1}{2}a_n-\dfrac{1}{2}(a_n-2a_{n+1})
  \]
  两侧分别取下极限和上极限,有
  \[
    \alpha=\dfrac{1}{2}(\alpha-A),
    \qquad
    \beta=\dfrac{1}{2}(\beta-A).
  \]
  因此 $\alpha=\beta=-A$,从而 $a_n\to-A$.
\end{solution}

%例题1.4.3
\begin{example}
  设 $a_n=nb_n$,其中 $\{b_n\}$ 满足 $|nb_n|\leq M$,且存在
  \[
    \lim\limits_{n\to\infty}n(b_n+b_{2n})=l,
  \]
  证明:$\{a_n\}$ 收敛,并求其极限.
\end{example}
\exampleblank

\begin{solution}
  由 $a_n=nb_n$,题设可写成
  \[
    a_n+\dfrac{1}{2}a_{2n}\longrightarrow l.
  \]
  又由 $|a_n|=|nb_n|\leq M$ 知 $\{a_n\}$ 有界,而上式等价于
  \[
    a_n-\dfrac{2l}{3}
    +\dfrac{1}{2}\left(a_{2n}-\dfrac{2l}{3}\right)\longrightarrow0.
  \]
  因此
  \[
    \varlimsup\limits_{n\to\infty}
    \left|a_n-\dfrac{2l}{3}\right|
    \leq\dfrac{1}{2}
    \varlimsup\limits_{n\to\infty}
    \left|a_{2n}-\dfrac{2l}{3}\right|
    \leq\dfrac{1}{2}
    \varlimsup\limits_{n\to\infty}
    \left|a_n-\dfrac{2l}{3}\right|.
  \]
  这个有限的非负上极限只能为 $0$,故 $a_n\to\dfrac{2l}{3}$.
\end{solution}

%例题1.4.4
\begin{example}
  设
  \[
    a_{n+1}=1-\dfrac{a_n^2}{b^2},
    \qquad 0<a_1<b,
    \qquad b>\sqrt2,
  \]
  讨论 $\{a_n\}$ 的敛散性.
\end{example}
\exampleblank

\begin{solution}
  由 $0<a_1<b$ 得 $0<a_2<1$.当 $0<a_n<1$ 时,因 $b>\sqrt2>1$,
  仍有 $0<a_{n+1}<1$.故从第二项起数列位于 $(0,1)$ 内.

  记
  \[
    \alpha=\varliminf\limits_{n\to\infty}a_n,
    \qquad
    \beta=\varlimsup\limits_{n\to\infty}a_n.
  \]
  函数 $1-\dfrac{x^2}{b^2}$ 在 $\lbrack0,1\rbrack$ 上连续且递减,所以
  \[
    \alpha=1-\dfrac{\beta^2}{b^2},
    \qquad
    \beta=1-\dfrac{\alpha^2}{b^2}.
  \]
  两式相减得到
  \[
    \beta-\alpha
    =\dfrac{(\beta-\alpha)(\alpha+\beta)}{b^2}.
  \]
  因 $0\leq\alpha\leq\beta\leq1$ 且 $b^2>2$,只能有 $\alpha=\beta$.
  故数列收敛,其极限是方程 $x=1-\dfrac{x^2}{b^2}$ 的正根
  \[
    \dfrac{b\sqrt{b^2+4}-b^2}{2}.
  \]
\end{solution}

%例题1.4.5
\begin{example}
  讨论下列数列 $\{a_n\}$ 的敛散性:
  \begin{enumerate}
    \item $\displaystyle a_{n+1}=\sqrt[3]{a_n+6}\quad(a_1>-6)$;
    \item $\displaystyle a_{n+1}=A\sqrt{a_n}+B\quad(A,B>0,\ a_1>0)$.
  \end{enumerate}
\end{example}
\exampleblank

\begin{solution}
  (1) $a_2=\sqrt[3]{a_1+6}>0$.取 $M=\max\{a_2,2\}$.因为
  \[
    0<\sqrt[3]{x+6}\leq M\qquad(0\leq x\leq M),
  \]
  其中右侧不等式来自 $M+6\leq M^3$,故 $0<a_n\leq M$ 对 $n\geq2$ 成立.
  记上下极限为 $\alpha,\beta$.
  由立方根函数连续且递增,
  \[
    \alpha=\sqrt[3]{\alpha+6},
    \qquad
    \beta=\sqrt[3]{\beta+6}.
  \]
  而 $x^3-x-6=(x-2)(x^2+2x+3)$,故 $\alpha=\beta=2$,即 $a_n\to2$.

  (2) 方程 $x=A\sqrt x+B$ 的唯一正根为
  \[
    L=\left(\dfrac{A+\sqrt{A^2+4B}}{2}\right)^2.
  \]
  对任意 $x>0$,
  \[
    A\sqrt x+B-x
    =(\sqrt L-\sqrt x)(\sqrt L+\sqrt x-A).
  \]
  因 $\sqrt L>A$,第二个因子恒正;又 $A\sqrt x+B$ 严格递增,所以递推始终留在
  $a_1$ 与 $L$ 之间,从而有界.记上下极限为 $\alpha,\beta$,由连续性和单调性,
  \[
    \alpha=A\sqrt\alpha+B,
    \qquad
    \beta=A\sqrt\beta+B.
  \]
  两者都是上述方程的唯一正根,故 $\alpha=\beta=L$.
\end{solution}

%例题1.4.6
\begin{example}
  设 $a_1>0$,
  \[
    a_{n+1}=\dfrac{1}{2}\left(a_n+\dfrac{1}{a_n}\right)+1,
  \]
  讨论 $\{a_n\}$ 敛散性.
\end{example}
\exampleblank

\begin{solution}
  由算术--几何平均值不等式,$a_2\geq2$,故 $a_n\in\lbrack2,+\infty)$ 对
  $n\geq2$ 成立.递推函数在该区间上的导数为
  $\dfrac{1}{2}-\dfrac{1}{2x^2}>0$,故严格递增;其不动点方程为
  $x^2-2x-1=0$,唯一的正不动点为
  $1+\sqrt2$.因此以 $a_2$ 和 $1+\sqrt2$ 为端点的闭区间在递推下保持不变,
  数列有界.

  记上下极限为 $\alpha,\beta$.由递推函数连续且递增,
  \[
    \alpha=\dfrac{1}{2}\left(\alpha+\dfrac{1}{\alpha}\right)+1,
    \qquad
    \beta=\dfrac{1}{2}\left(\beta+\dfrac{1}{\beta}\right)+1.
  \]
  由于该方程只有一个正根 $1+\sqrt2$,故 $\alpha=\beta=1+\sqrt2$.
\end{solution}

%例题1.4.7
\begin{example}
  设
  \[
    y_n=px_n+qx_{n+1},
    \qquad |p|<|q|,
  \]
  证明:$\{x_n\}$ 与 $\{y_n\}$ 同敛散.
\end{example}
\exampleblank

\begin{solution}
  若 $x_n\to X$,则立即有
  \[
    y_n=px_n+qx_{n+1}\longrightarrow(p+q)X.
  \]

    反之,设 $y_n\to Y$.由于 $|p|<|q|$,有 $q\neq0$ 且 $p+q\neq0$.
  由
  \[
    x_{n+1}=-\dfrac{p}{q}x_n+\dfrac{y_n}{q},
  \]
  逐次代入得
  \[
    x_2=-\dfrac{p}{q}x_1+\dfrac{y_1}{q},
  \]
  \[
    x_3=
    \left(-\dfrac{p}{q}\right)^2x_1
    +\dfrac{1}{q}
    \left[
      -\dfrac{p}{q}y_1+y_2
    \right],
  \]
  继续迭代可得
  \[
    x_n=
    \left(-\dfrac{p}{q}\right)^{n-1}x_1
    +\dfrac{1}{q}
    \sum\limits_{k=1}^{n-1}
    \left(-\dfrac{p}{q}\right)^{n-1-k}y_k.
  \]
  因 $y_n\to Y$,故 $\{y_n\}$ 有界.于是
  \[
    \begin{aligned}
      |x_n|
      &\leq
      \left|\dfrac{p}{q}\right|^{n-1}|x_1|
      +\dfrac{\displaystyle\sup\limits_k|y_k|}{|q|}
      \sum\limits_{k=1}^{n-1}
      \left|\dfrac{p}{q}\right|^{n-1-k} \\
      &=
      \left|\dfrac{p}{q}\right|^{n-1}|x_1|
      +\dfrac{\displaystyle\sup\limits_k|y_k|}{|q|}
      \sum\limits_{j=0}^{n-2}
      \left|\dfrac{p}{q}\right|^j.
    \end{aligned}
  \]
   由于 $\left|\dfrac{p}{q}\right|<1$,右端一致有界,故 $\{x_n\}$ 有界.
  记
  \[
    A=\varlimsup\limits_{n\to\infty}x_n,
    \qquad
    B=\varliminf\limits_{n\to\infty}x_n,
    \qquad
    r=-\dfrac{p}{q}.
  \]
  则 $|r|<1$,且
  \[
    x_{n+1}=rx_n+\dfrac{y_n}{q},
    \qquad
    \dfrac{y_n}{q}\longrightarrow\dfrac{Y}{q}.
  \]

  若 $r\geq0$,则
  \[
    A=rA+\dfrac{Y}{q},
    \qquad
    B=rB+\dfrac{Y}{q},
  \]
  从而 $A=B$.

  若 $r<0$,则
  \[
    A=rB+\dfrac{Y}{q},
    \qquad
    B=rA+\dfrac{Y}{q}.
  \]
  两式相减得
  \[
    (1+r)(A-B)=0.
  \]
  由于 $|r|<1$,故 $A=B$.

  因此 $\{x_n\}$ 收敛.设 $x_n\to X$,则
  \[
    Y=(p+q)X,
  \]
  即
  \[
    X=\dfrac{Y}{p+q}.
  \]
\end{solution}


%例题1.4.8
\begin{example}

  证明下列命题:

  \begin{enumerate}

    \item 若非负数列 $\{a_n\}$ 满足对任意正整数 $m,n$ 都有
          \[
            a_{m+n}\leq a_m+a_n,
          \]
          则数列 $\left\{\dfrac{a_n}{n}\right\}$ 收敛;

    \item 若非负数列 $\{a_n\}$ 满足对任意正整数 $m,n$ 都有
          \[
            a_{m+n}\leq a_m a_n,
          \]
          则数列 $\{\sqrt[n]{a_n}\}$ 收敛.

  \end{enumerate}

\end{example}

\exampleblank

\begin{solution}

  (1) 令
  \[
    \alpha=\inf\limits_{k\geq 1}\dfrac{a_k}{k}.
  \]
  显然
  \[
    \varliminf\limits_{n\to\infty}\dfrac{a_n}{n}\geq\alpha.
  \]

  固定正整数 $k$.将 $n$ 写成
  \[
    n=qk+r,\qquad 0\leq r<k.
  \]
  由题设反复使用次可加性,得到
  \[
    a_n\leq qa_k+a_r,
  \]
  其中 $r=0$ 时略去末项.

  令
  \[
    C_k=\max\{a_1,\ldots,a_{k-1}\},
  \]
  则
  \[
    \dfrac{a_n}{n}
    \leq
    \dfrac{qa_k+C_k}{qk+r}.
  \]
  令 $n\to\infty$,即 $q\to\infty$,可得
  \[
    \varlimsup\limits_{n\to\infty}\dfrac{a_n}{n}
    \leq
    \dfrac{a_k}{k}.
  \]
  由于 $k$ 任意,对 $k$ 取下确界,得到
  \[
    \varlimsup\limits_{n\to\infty}\dfrac{a_n}{n}
    \leq\alpha.
  \]
  因此
  \[
    \varlimsup\limits_{n\to\infty}\dfrac{a_n}{n}
    \leq
    \alpha
    \leq
    \varliminf\limits_{n\to\infty}\dfrac{a_n}{n},
  \]
  从而
  \[
    \lim\limits_{n\to\infty}\dfrac{a_n}{n}
    =\alpha.
  \]

  (2) 若存在正整数 $k$ 使 $a_k=0$,则对任意 $n>k$,
  \[
    0\leq a_n\leq a_k a_{n-k}=0,
  \]
  故 $a_n=0$ 对所有 $n\geq k$ 成立,于是
  \[
    \lim\limits_{n\to\infty}\sqrt[n]{a_n}=0.
  \]

  以下设 $a_n>0$ 对所有正整数 $n$ 成立.令
  \[
    b_n=\ln a_n.
  \]
  由
  \[
    a_{m+n}\leq a_m a_n
  \]
  得
  \[
    b_{m+n}\leq b_m+b_n.
  \]
  因此由第 (1) 问,
  \[
    \lim\limits_{n\to\infty}\dfrac{b_n}{n}
  \]
  存在.

  又因为
  \[
    \sqrt[n]{a_n}
    =a_n^{\frac{1}{n}}
    =\exp\left(\dfrac{b_n}{n}\right),
  \]
  且指数函数连续,所以 $\{\sqrt[n]{a_n}\}$ 收敛.

\end{solution}


%例题1.4.9
\begin{example}

  证明下列命题:

  \begin{enumerate}

    \item 设 $\{a_n\}$ 是正数列,则
          \[
            \varlimsup\limits_{n\to\infty}
            n\left(\dfrac{1+a_{n+1}}{a_n}-1\right)
            \geq1;
          \]

    \item 设 $\{a_n\}$ 是正数列,则
          \[
            \varlimsup\limits_{n\to\infty}
            \left(\dfrac{a_1+a_{n+1}}{a_n}\right)^n
            \geq e.
          \]

  \end{enumerate}

\end{example}

\exampleblank

\begin{solution}

  (1) 若结论不成立,则可取 $q<1$,使得当 $n$ 充分大时
  \[
    n\left(\dfrac{1+a_{n+1}}{a_n}-1\right)\leq q.
  \]
  因而
  \[
    \dfrac{a_{n+1}}{n+1}
    \leq\dfrac{n+q}{n+1}\dfrac{a_n}{n}-\dfrac{1}{n+1}
    <\dfrac{a_n}{n}-\dfrac{1}{n+1}.
  \]
  从某个固定的 $N$ 起迭代此式,得到
  \[
    \dfrac{a_{m+1}}{m+1}
    <\dfrac{a_N}{N}-\sum\limits_{n=N}^m\dfrac{1}{n+1}.
  \]
  右端最终为负,与 $a_{m+1}>0$ 矛盾,故 (1) 成立.

  (2) 对正数列 $\dfrac{a_n}{a_1}$ 应用 (1),得到
  \[
    \varlimsup\limits_{n\to\infty}
    n\left(\dfrac{a_1+a_{n+1}}{a_n}-1\right)\geq1.
  \]
  若 (2) 左端小于 $e$,则存在 $c<1$,使得当 $n$ 充分大时
  \[
    \left(\dfrac{a_1+a_{n+1}}{a_n}\right)^n\leq e^c.
  \]
  于是
  \[
    n\left(\dfrac{a_1+a_{n+1}}{a_n}-1\right)
    \leq n\left(e^{\frac{c}{n}}-1\right)\longrightarrow c<1,
  \]
  与前式矛盾.故 (2) 成立.

\end{solution}
